\partchapter{2}{第二部分\quad 系统组成及功能说明}

\Needspace{85pt}

\section{整体介绍}

处理器采用经典五级顺序流水线结构：取指（IF）、译码与读寄存器（ID）、执行（EX）、访存（MEM）、写回（WB），并配置 I-cache、D-cache、分支预测器、旁路网络与集中式冒险控制。顶层 \texttt{cpu\_\allowbreak{}core} 之下的并列单元构成如下表。

{\small
\renewcommand{\thetable}{2-1}
\begin{longtable}{@{}>{\raggedright\arraybackslash}p{\dimexpr 0.0700\linewidth-2\tabcolsep\relax}>{\raggedright\arraybackslash}p{\dimexpr 0.4650\linewidth-2\tabcolsep\relax}>{\raggedright\arraybackslash}p{\dimexpr 0.4650\linewidth-2\tabcolsep\relax}@{}}
\caption{处理器功能单元组成}\label{tab:2-1}\\
\toprule
\textbf{单元} & \textbf{文件} & \textbf{职责} \\
\midrule
\endfirsthead
\multicolumn{3}{@{}l}{\footnotesize 续表 2-1}\\
\toprule
\textbf{单元} & \textbf{文件} & \textbf{职责} \\
\midrule
\endhead
取指 & \texttt{if\_\allowbreak{}\allowbreak{}stage.\allowbreak{}v}、\texttt{if\_\allowbreak{}\allowbreak{}id.\allowbreak{}v} & PC 生成、I-cache 查询、预测目标选择、冲刷响应 \\
译码 & \texttt{id\_\allowbreak{}\allowbreak{}stage.\allowbreak{}v}、\texttt{id\_\allowbreak{}\allowbreak{}ex.\allowbreak{}v} & 指令字段解析、立即数生成、结果类型（ResultKind）判定、寄存器索引下发 \\
执行 & \texttt{ex\_\allowbreak{}\allowbreak{}stage.\allowbreak{}v}、\texttt{mult.\allowbreak{}v}、\texttt{divider.\allowbreak{}v} & ALU 与移位、分支条件判定、多周期乘除、Ascon ISE 接入 \\
访存 & \texttt{mem\_\allowbreak{}\allowbreak{}stage.\allowbreak{}v}、\texttt{mem\_\allowbreak{}\allowbreak{}wb.\allowbreak{}v} & load/store 对齐与 D-cache 事务、总线请求 \\
写回 & \texttt{mem\_\allowbreak{}\allowbreak{}wb.\allowbreak{}v}、\texttt{regfile.\allowbreak{}v} & 结果选择、GPR 写端口、CSR 写端口 \\
控制 & \texttt{ctrl.\allowbreak{}v}（\texttt{stall\_\allowbreak{}\allowbreak{}ctrl}）、\texttt{bypass.\allowbreak{}v} & 5 位 \texttt{stall/\allowbreak{}flush} 向量、两级旁路选择 \\
存储 & \texttt{icache.\allowbreak{}v}、\texttt{dcache.\allowbreak{}v}、\texttt{bus\_\allowbreak{}\allowbreak{}interconnect.\allowbreak{}v} & 各 64 组 $\times$ 4 字/行；地址译码与总线状态机 \\
系统 & \texttt{csr\_\allowbreak{}\allowbreak{}file.\allowbreak{}v}、\texttt{clint.\allowbreak{}v}、\texttt{uart\_\allowbreak{}\allowbreak{}ctrl.\allowbreak{}v}、\texttt{gpio.\allowbreak{}v}、\texttt{cpi\_\allowbreak{}\allowbreak{}cnt.\allowbreak{}v} & 特权 CSR、定时器中断、串口、GPIO、CPI 事件计数 \\
互连 & \texttt{cpu\_\allowbreak{}\allowbreak{}axi\_\allowbreak{}\allowbreak{}bridge.\allowbreak{}v}、\texttt{axi\_\allowbreak{}\allowbreak{}full\_\allowbreak{}\allowbreak{}master.\allowbreak{}v}、\texttt{soc\_\allowbreak{}\allowbreak{}top.\allowbreak{}v}、\texttt{top.\allowbreak{}v}、\texttt{top\_\allowbreak{}\allowbreak{}pl\_\allowbreak{}\allowbreak{}bram.\allowbreak{}v} & 内部总线至 AXI4 master、BRAM 与外设；板级顶层与时钟复位 \\
\bottomrule
\end{longtable}
}

\figpdf{2-1}{fig_2_1}{处理器与 SoC 总体框图}{CPU 五级流水、I/D-cache、旁路与 冒险控制、总线互连、AXI BRAM、UART/GPIO/CLINT/CPI，以及 Ascon MMIO 与 EX 级 ISE。 当前系统时钟由 PS7 FCLK0 提供，PL-only 顶层仅作历史对照。}

数据通路的关键接口约定如下：

\begin{itemize}
\tightlist
\item
  \textbf{指令与数据的地址空间由 Cache 区分}：I-cache 与 D-cache 独立查询同一片 BRAM 区域（统一编址， \textgreater{} \texttt{0x1000\_\allowbreak{}0000} 起 128 KiB）；不可 Cache 的访问（MMIO）直接经总线完成，不进入 Cache；
\item
  \textbf{结果类型于 ID 级判定}：指令在译码阶段即标注结果来源（算术、逻辑、移位、load、跳转链接、乘法）， EX 级据此选择数据通路，避免在执行级重复译码；
\item
  \textbf{旁路于 ID 级选路、于 EX 入口生效}：优先级为 \texttt{EX/MEM\ \textgreater{}\ MEM/WB}（\texttt{bypass.\allowbreak{}v} 中两级 MUX）； load-use 冒险以停顿方式解决，不采用推测加重放；
\item
  \textbf{CSR 读写置于 EX 级}：\texttt{mepc}、\texttt{mcause}、\texttt{mtvec} 与陷入入口位于同一路径，使 \texttt{ecall} 与 \texttt{mret} 的 PC 修复仅有单一实现。
\end{itemize}

\Needspace{85pt}

\section{处理器架构与设计思路}

\Needspace{85pt}

\subsection{设计目标与总体思路}

处理器以 CICC1000023 作品中的自研 RV32 五级流水 CPU 为基础，本作品不重新设计 CPU，而是在其上完成器件迁移、时序改造与密码加速扩展。设计围绕三个目标展开：

\begin{enumerate}
\def\labelenumi{\arabic{enumi}.}
\tightlist
\item
  \textbf{在 XC7Z020 上以 100 MHz 稳定运行}。7020 的 LUT 与布线资源有限，结构选择优先考虑时序可收敛，而不是追求更高的单线程性能。
\item
  \textbf{为三种 Ascon 实现提供公平的对照平台}。软件、MMIO 与定制指令必须运行在同一 CPU、同一存储层次、同一测量方式（\texttt{mcycle}）下，周期差异才只来自实现方式本身。
\item
  \textbf{让定制指令接入后不破坏既有流水线语义}。新增指令要能被停顿、被陷入取消，且不会因停顿而重复执行。
\end{enumerate}

据此形成以下设计取舍：

{\small
\renewcommand{\thetable}{2-2}
\begin{longtable}{@{}>{\raggedright\arraybackslash}p{\dimexpr 0.1248\linewidth-2\tabcolsep\relax}>{\raggedright\arraybackslash}p{\dimexpr 0.4376\linewidth-2\tabcolsep\relax}>{\raggedright\arraybackslash}p{\dimexpr 0.4376\linewidth-2\tabcolsep\relax}@{}}
\caption{处理器设计取舍}\label{tab:2-2}\\
\toprule
\textbf{决策} & \textbf{选择} & \textbf{理由与代价} \\
\midrule
\endfirsthead
\multicolumn{3}{@{}l}{\footnotesize 续表 2-2}\\
\toprule
\textbf{决策} & \textbf{选择} & \textbf{理由与代价} \\
\midrule
\endhead
流水线深度 & 经典五级顺序流水，不做乱序与多发射 & 控制逻辑可审查、可验证；单发射使周期计数可直接比较 \\
冒险控制 & \texttt{ctrl.\allowbreak{}v} 集中产生 stall / flush，按固定优先级裁决 & 多种停顿源（Cache、总线、乘除、ISE、load-use）叠加时只有一处裁决，便于证明正确性 \\
数据冒险 & EX/MEM、MEM/WB 两级旁路，load-use 停顿一拍 & 不做推测加重放，以少量停顿换取简单可靠的控制 \\
译码 & ID 级标注结果类型（ResultKind）并随流水传递 & 后级按类型选数据通路，避免在 EX/WB 重复译码，缩短关键路径 \\
存储层次 & 分离 I/D-cache，代码与数据统一编址于 128 KiB BRAM；MMIO 绕过 Cache & 取指与访存并行；外设访问保持顺序与副作用语义 \\
分支预测 & 64 项 BTB + 2-bit BHT，分布式 RAM 零延迟读出 & 小表面积与时序代价低；命中率未单独实测，不作性能结论 \\
加速接入 & MMIO 外设挂总线；ISE 置于 EX 级、持有 320-bit 私有状态 & 前者不改 CPU，后者省去总线往返；ISE 以 hold / ack 握手跨多拍执行 \\
\bottomrule
\end{longtable}
}

时序收敛同样体现上述思路：先用 33 组实现策略确定工具能回收的余量，再按失败端点归因做结构改造（D-cache 增加 MISS\_CAPTURE、乘法器寄存器化、ISE 局部复位），而不是降低频率（见 2.2.11 节）。

\Needspace{85pt}

\subsection{流水线微架构}

图 2-2 按流水级分列给出处理器微架构。顶部为贯穿各级的集中冒险控制，竖条为四组流水寄存器，底部为总线与 SoC 外设。

\figpdf{2-2}{fig_2_12}{RV32 处理器流水线微架构}{五级流水各级功能单元、流水寄存器、 集中冒险控制、旁路与写回、分支重定向与预测器训练，以及 I/D-cache 与 MMIO 的总线出口。}

\begin{itemize}
\tightlist
\item
  \textbf{IF}：PC 按''陷入向量 \textgreater{} 分支重定向 \textgreater{} 预测目标 \textgreater{} PC+4''选择；分支预测器在同拍组合读出目标， I-cache 命中当拍返回指令，未命中则停顿取指并整行回填。
\item
  \textbf{ID}：并行生成五种立即数，标注结果类型；寄存器堆双读一写并带写读直通； load-use 冒险在此检测，命中时冻结前两级并插入气泡。
\item
  \textbf{EX}：旁路网络先选出最新操作数，再送入 ALU、乘除法、CSR 与 Ascon ISE。分支在 EX 级裁决，误预测时重定向 PC 并冲刷两条较年轻指令，同时训练 BTB/BHT。乘除与 ISE 为多周期单元，通过 \texttt{ex\_\allowbreak{}stall\_\allowbreak{}req} 冻结前级。
\item
  \textbf{MEM}：可 Cache 地址走两路组相联 D-cache；MMIO 地址绕过 Cache，直接发出单拍注册请求。
\item
  \textbf{WB}：按结果类型选择 ALU、load、PC+4 或 CSR 值写回寄存器堆，并经旁路回送 EX。
\end{itemize}

各单元的实现细节见 2.2.3--2.2.10 节，结构改造汇总见 2.2.11 节。

\Needspace{85pt}

\subsection{取指单元（IF）}\label{ux53d6ux6307ux5355ux5143if}

取指级的职责为每周期产生一个 PC 及一条指令，需同时处理 Cache 命中、分支预测与流水线冲刷三项相互耦合的条件。

\begin{itemize}
\tightlist
\item
  \textbf{PC 来源优先级}（由高至低）：陷入与中断入口、EX 级分支重定向（\texttt{jump\_\allowbreak{}en}）、分支预测目标、 \texttt{pc\ +\ 4}；
\item
  \textbf{I-cache 查询}：容量 1 KiB，64 组 $\times$ 4 字/行，直接映射；tag 命中时当拍返回。未命中时进入 \texttt{FILL} 状态并向总线发出一次 burst 请求，\texttt{ctrl.\allowbreak{}v} 收到 \texttt{icache\_\allowbreak{}stall} 后冻结 IF 与 IF/ID；
\item
  \textbf{预测器接口}：取指级以当前 PC 组合读出 \texttt{pred\_\allowbreak{}taken} 与 \texttt{pred\_\allowbreak{}target}（分布式 RAM，零延迟），预测命中时当拍改选 PC，不增加流水级；
\item
  \textbf{冲刷}：\texttt{flush{[}4{]}} 清除 IF，\texttt{flush{[}3{]}} 清除 IF/ID。分支重定向（\texttt{jump\_\allowbreak{}en}）时 \texttt{ctrl.\allowbreak{}v} 置 \texttt{flush{[}3:2{]}}（清除 IF/ID 与 ID/EX 两级，作废两条较年轻指令）；陷入、\texttt{mret} 与外部中断时置 \texttt{flush{[}3:1{]}}（额外清除 EX/MEM）。各级冲刷的流水线代价随清除级数而变，本文不作未经实测的拍数结论。
\end{itemize}

\figpdf{2-3}{fig_2_3}{IF 级数据通路与冲刷语义}{PC 选择、I-cache 查询、 预测改选、miss FILL 与分支冲刷（flush[3:2] 两级）；为结构示意，不作为实测周期波形。}

\Needspace{85pt}

\subsection{译码单元（ID）}\label{ux8bd1ux7801ux5355ux5143id}

译码级完成字段解析与指令意图标注，为流水线中控制逻辑最密集的一级。

\begin{itemize}
\tightlist
\item
  \textbf{立即数生成}：I、S、B、U、J 五种格式并行生成，由 opcode 选择；
\item
  \textbf{寄存器堆}：32 $\times$ 32-bit，双读一写；\texttt{x0} 恒为零（写端口屏蔽）；同步写、组合读。读地址与有效写地址相同时直接返回 \texttt{w\_\allowbreak{}reg\_\allowbreak{}data}，存在寄存器堆内部写读旁路，并非返回旧值；
\item
  \textbf{结果类型（ResultKind）}：将指令映射为六类结果来源（算术、逻辑、移位、load、跳转链接、乘法），该标注随流水线传递至 EX/MEM/WB，为后级''不再重复译码''的基础；
\item
  \textbf{CSR 指令与陷入}：\texttt{csrrw}、\texttt{csrrs}、\texttt{csrrc}、\texttt{ecall} 与 \texttt{mret} 于 ID 级识别为 SYSTEM 类，具体读写与 \texttt{mepc}、\texttt{mcause} 更新在 EX 级完成；
\item
  \textbf{load-use 检测}：ID 级比较 \texttt{id\_\allowbreak{}ex.\allowbreak{}rd} 与当前 \texttt{rs1}、\texttt{rs2}，命中时置 \texttt{stall\_\allowbreak{}ld}，由 \texttt{ctrl.\allowbreak{}v} 冻结 IF 与 IF/ID（\texttt{stall{[}4:3{]}}），并在 ID/EX 插入气泡（\texttt{flush{[}2{]}}）；EX/MEM 不再清除。
\end{itemize}

\Needspace{85pt}

\subsection{执行单元（EX）}

执行级为多功能并行与多周期串行相结合的混合结构，本作品的两处关键改造均作用于该级。

\textbf{并行功能单元}：ALU（加、减、逻辑、比较、移位）、分支条件判定、跳转链接地址生成与 CSR 读写均在单周期内完成；乘法器 \texttt{mult.\allowbreak{}v} 完成 32 $\times$ 32 有符号与无符号乘法（RV32M 的 \texttt{mul}、\texttt{mulh}、\texttt{mulhsu}、\texttt{mulhu}），综合后被映射为 4 个 DSP48E1（实现报告 DSP 利用率 4／220），映射结果由工具的乘法推断产生而非显式例化 DSP 原语；除法器 \texttt{divider.\allowbreak{}v} 为迭代式多周期实现（\texttt{div}、\texttt{divu}、\texttt{rem}、\texttt{remu}）。

\textbf{多周期与握手}：执行级向 \texttt{ctrl.\allowbreak{}v} 发出 \texttt{ex\_\allowbreak{}stall\_\allowbreak{}req}，后者冻结 IF、IF/ID 与 ID/EX 并在 EX/MEM 插入气泡。多周期单元完成时将结果与 \texttt{done} 一并锁存，直至被确认（ack）方予释放。乘法器于本轮改造中增加 \texttt{mul\_\allowbreak{}ack} 与 \texttt{cancel} 两个内部端口，其契约为：输入寄存器至乘积寄存器，接受后下一拍乘积与 \texttt{mul\_\allowbreak{}done} 有效；\texttt{done} 与结果保持至 ack；cancel 或 reset 取消且不产生完成信号。该契约使''乘法结果在 EX 级保持''与''陷入取消''两项行为互不干扰。

\textbf{Ascon ISE 接入点}：custom-0（opcode \texttt{0x0b}，funct7 = 0，funct3 区分四类操作）在该级被识别为一类带状态的多周期指令，其读写对象为 ISE 内部的 10 个 32-bit 状态字（§2.3.3.4），而非通用寄存器。执行级以 hold 方式使流水线跨越 P12/P8 的多拍执行，\texttt{done} 保持至指令推进；``完成有效方予退休''的判据保证该指令不会因跨过完成的停顿而被重复执行。

\figpdf{2-4}{fig_2_4}{EX 级功能单元与多周期握手状态图}{含 `Idle $\rightarrow$ Busy $\rightarrow$ Done $\rightarrow$ ack 转换、陷入与中断取消分支，以及 ISE hold 与 ALU 单拍路径的复用关系。}

\Needspace{85pt}

\subsection{访存与写回（MEM/WB）及 Cache 结构}\label{ux8bbfux5b58ux4e0eux5199ux56dememwbux53ca-cache-ux7ed3ux6784}

访存级将执行级的地址与数据转换为 D-cache 事务或总线事务；写回级完成 GPR 与 CSR 写回。

\textbf{D-cache 结构}（\texttt{dcache.\allowbreak{}v}）：容量 2 KiB，64 组 $\times$ 4 字/行，两路组相联，采用 write-through 策略（store 不等待填充），tag 与数据分存；不可 Cache 的地址（MMIO 窗口）绕过 Cache，改为单拍注册请求。状态机为 \texttt{IDLE\ /\allowbreak{}\ MISS\_\allowbreak{}CAPTURE\ /\allowbreak{}\ FILL} 三态：

\begin{itemize}
\tightlist
\item
  \texttt{IDLE}：查询 tag，命中即返回；
\item
  \texttt{MISS\_\allowbreak{}CAPTURE}：本轮新增状态，将 miss 的 tag、index、way 及请求属性锁存一拍，使''miss 判定至总线请求''不再经过单一长组合路径；
\item
  \texttt{FILL}：等待总线填充完成，逐字写入 way。
\end{itemize}

\textbf{I-cache}：容量 1 KiB，64 组 $\times$ 4 字/行，直接映射，状态机为 \texttt{IDLE/\allowbreak{}FILL} 两态，未命中时整行填充。

\textbf{写回}：\texttt{mem\_\allowbreak{}wb} 选择访存结果、load 数据、PC + 4 或 CSR 读值，驱动 \texttt{regfile} 写口；\texttt{x0} 于写口屏蔽， ``读 x0 恒为零''由寄存器堆本身保证，写回级不再作特判。

\figpdf{2-5}{fig_2_5}{D-cache 与总线请求时序}{命中单拍、`MISS\_CAPTURE` 一拍锁存、 `FILL` 逐字写入，以及 store 与 MMIO 的绕过通路。}

\Needspace{85pt}

\subsection{CSR、异常与计数}\label{csrux5f02ux5e38ux4e0eux8ba1ux6570}

\texttt{csr\_\allowbreak{}file.\allowbreak{}v} 实现的 CSR 集合为：\texttt{mstatus}（0x300）、\texttt{misa}（0x301）、\texttt{mie}（0x304）、\texttt{mtvec}（0x305）、 \texttt{mscratch}（0x340）、\texttt{mepc}（0x341）、\texttt{mcause}（0x342）、\texttt{mtval}（0x343）、\texttt{mip}（0x344）、 \texttt{mvendorid}（0xF11）、\texttt{marchid}（0xF12）、\texttt{mimpid}（0xF13）、\texttt{mhartid}（0xF14），以及 64-bit \texttt{mcycle} 与 \texttt{mcycleh}（0xB00／0xB80）。其中 \texttt{mvendorid}／\texttt{marchid}／\texttt{mimpid}／\texttt{mhartid} 返回常量 0。

\begin{itemize}
\tightlist
\item
  \texttt{mcycle} 为\textbf{唯一周期来源}：每个非复位时钟加一，RV32 经 \texttt{0xb00}（低）与 \texttt{0xb80}（高）访问，软件按 high$\rightarrow$low$\rightarrow$high 顺序读取并校验一致性后重读。本作品全部周期加速比均以 \texttt{mcycle} 差值计算，不以事件计数之和替代------后者（\texttt{cpi\_\allowbreak{}cnt} 的 inst、stall、flush 事件）存在重叠，不可相加为周期数；
\item
  \textbf{陷入流程}：\texttt{ecall} 于 EX 级写入 \texttt{mepc}（当前 PC）与 \texttt{mcause\ =\ 11}，跳转至 \texttt{mtvec}；\texttt{mret} 自 \texttt{mepc} 恢复 PC；外部中断经 \texttt{ireq\_external\ =\ gpio\_irq\ \textbar{}\ uart\_irq\ \textbar{}\ ascon\_irq} 汇聚，抢占时由 \texttt{stall\_\allowbreak{}ctrl} 清除较年轻指令；
\item
  \textbf{计数与 CPI}：\texttt{cpi\_\allowbreak{}cnt.\allowbreak{}v} 挂接 \texttt{0x6000\_\allowbreak{}0000} 窗口，提供 \texttt{probe\_\allowbreak{}inst\_\allowbreak{}valid}、\texttt{probe\_\allowbreak{}stall\_\allowbreak{}ld}、 \texttt{probe\_\allowbreak{}stall\_\allowbreak{}md} 与 \texttt{probe\_\allowbreak{}flush\_\allowbreak{}br}（每事件计权 2）四组 64-bit 计数；CPI 以同一区间的 cycle delta 与已退休指令 delta 计算，不采用事件加权和。
\end{itemize}

\Needspace{85pt}

\subsection{分支预测}

预测器为 64 项单表结构（\texttt{branch\_\allowbreak{}predictor.\allowbreak{}v}，\texttt{INDEX\_\allowbreak{}WIDTH\ =\ 6}）：

\begin{itemize}
\tightlist
\item
  \textbf{表项}：tag 表、BTB 目标表（分布式 RAM，\texttt{ram\_\allowbreak{}style="distributed"}，整体 \texttt{dont\_\allowbreak{}touch}）、2-bit 饱和计数器 BHT 与打包 valid 向量；
\item
  \textbf{索引}：\texttt{index\ =\ pc{[}7:2{]}}（指令 4 字节对齐，低 2 位恒为零），\texttt{tag\ =\ pc{[}31:8{]}}；
\item
  \textbf{预测}：仅当 \texttt{valid\ \&\&\ tag\_match\ \&\&\ bht\ \textgreater{}=\ 2\textquotesingle{}b10} 时判为 taken，目标取 BTB 中记录的历史目标；
\item
  \textbf{更新}：分支进入 EX 级后训练（\texttt{update\_\allowbreak{}en}、\texttt{pc\_\allowbreak{}ex}、\texttt{actual\_\allowbreak{}taken}、\texttt{actual\_\allowbreak{}target}），并按 2-bit 状态机更新 BHT（\texttt{00↔01↔10↔11}，初值为弱不跳）。误预测冲刷由 \texttt{ctrl.\allowbreak{}v} 统一处理，预测器不参与冲刷。
\end{itemize}

采用''小型表 + 分布式 RAM''而非大容量 BTB，系针对本平台的容量与时序折衷：在 100 MHz、LUT 资源受限的 7020 器件上，该结构的面积与时序代价均较低；就本 SoC 的主要负载（AEAD 分块循环与基准测试主循环）而言，其命中率与单独性能贡献尚未量化，本文不以结构描述替代实测结论。

\Needspace{85pt}

\subsection{总线与 SoC 外设}\label{ux603bux7ebfux4e0e-soc-ux5916ux8bbe}

\texttt{bus\_\allowbreak{}interconnect.\allowbreak{}v} 按地址高位将读写请求译码至各从设备；\texttt{cpu\_\allowbreak{}axi\_\allowbreak{}bridge.\allowbreak{}v} 与 \texttt{axi\_\allowbreak{}full\_\allowbreak{}master.\allowbreak{}v} 将内部事务转换为 AXI4 master（ID 1 bit、WSTRB = XLEN/8，支持单拍与 burst）。

{\small
\renewcommand{\thetable}{2-3}
\begin{longtable}{@{}>{\raggedright\arraybackslash}p{\dimexpr 0.0832\linewidth-2\tabcolsep\relax}>{\raggedright\arraybackslash}p{\dimexpr 0.1445\linewidth-2\tabcolsep\relax}>{\raggedright\arraybackslash}p{\dimexpr 0.2984\linewidth-2\tabcolsep\relax}>{\raggedright\arraybackslash}p{\dimexpr 0.4738\linewidth-2\tabcolsep\relax}@{}}
\caption{SoC 地址窗口与外设}\label{tab:2-3}\\
\toprule
\textbf{窗口} & \textbf{基址} & \textbf{从设备} & \textbf{说明} \\
\midrule
\endfirsthead
\multicolumn{4}{@{}l}{\footnotesize 续表 2-3}\\
\toprule
\textbf{窗口} & \textbf{基址} & \textbf{从设备} & \textbf{说明} \\
\midrule
\endhead
MEM & \texttt{0x1000\_\allowbreak{}\allowbreak{}0000} & AXI BRAM Controller $\rightarrow$ Block RAM & 128 KiB，代码与数据统一空间，可 Cache \\
CLINT & \texttt{0x1200\_\allowbreak{}\allowbreak{}0000} & \texttt{clint.\allowbreak{}v} & \texttt{msip}、\texttt{mtimecmp}、\texttt{mtime}，提供软件定时中断 \\
UART & \texttt{0x2000\_\allowbreak{}\allowbreak{}0000} & \texttt{uart\_\allowbreak{}\allowbreak{}ctrl.\allowbreak{}v} & 当前 PYNQ 内部链路 6.25 Mbaud；历史板外串口 115200；读 FIFO 有副作用 \\
GPIO & \texttt{0x2001\_\allowbreak{}\allowbreak{}0000} & \texttt{gpio.\allowbreak{}v} & LED 控制、按键事件锁存（读清），两级同步与消抖 \\
ASCON & \texttt{0x3000\_\allowbreak{}\allowbreak{}0000} & \texttt{rtl\_\allowbreak{}\allowbreak{}p0/\allowbreak{}ascon\_\allowbreak{}\allowbreak{}mmio.\allowbreak{}v} & AEAD 寄存器接口，由调用者写 key（§2.3.3.3） \\
CPI & \texttt{0x6000\_\allowbreak{}\allowbreak{}0000} & \texttt{cpi\_\allowbreak{}\allowbreak{}cnt.\allowbreak{}v} & 四组 64-bit 事件计数 \\
\bottomrule
\end{longtable}
}

MMIO 事务不进入 D-cache。当前发布工程显式选择 \texttt{rtl\_\allowbreak{}p0/\allowbreak{}soc\_\allowbreak{}top.\allowbreak{}v} 与 \texttt{rtl\_\allowbreak{}p0/\allowbreak{}ascon\_\allowbreak{}mmio.\allowbreak{}v}；不能直接用工作树中同名但不同版本的文件替代发布源集。

\Needspace{85pt}

\subsection{流水线控制与冒险处理}

\texttt{ctrl.\allowbreak{}v} 输出两个 5 位向量 \texttt{stall{[}4:0{]}} 与 \texttt{flush{[}4:0{]}}，位序为 \texttt{\{IF,\ IF/ID,\ ID/EX,\ EX/MEM,\ MEM/WB\}}，优先级由高至低如下表。

{\small
\renewcommand{\thetable}{2-4}
\begin{longtable}{@{}>{\raggedright\arraybackslash}p{\dimexpr 0.0744\linewidth-2\tabcolsep\relax}>{\raggedright\arraybackslash}p{\dimexpr 0.4836\linewidth-2\tabcolsep\relax}>{\raggedright\arraybackslash}p{\dimexpr 0.4420\linewidth-2\tabcolsep\relax}@{}}
\caption{流水线控制事件优先级}\label{tab:2-4}\\
\toprule
\textbf{顺序} & \textbf{事件} & \textbf{\texttt{ctrl.\allowbreak{}v} 中的动作} \\
\midrule
\endfirsthead
\multicolumn{3}{@{}l}{\footnotesize 续表 2-4}\\
\toprule
\textbf{顺序} & \textbf{事件} & \textbf{\texttt{ctrl.\allowbreak{}v} 中的动作} \\
\midrule
\endhead
1 & \texttt{inter\_\allowbreak{}\allowbreak{}req}（外部中断）或 \texttt{jump\_\allowbreak{}\allowbreak{}trap\_\allowbreak{}\allowbreak{}en}（陷入，同一条件分支） & \texttt{flush[3:1]}；\texttt{jump\_\allowbreak{}\allowbreak{}sel = 10}（中断／陷入向量） \\
2 & \texttt{jump\_\allowbreak{}\allowbreak{}trap\_\allowbreak{}\allowbreak{}out}（\texttt{mret}） & \texttt{flush[3:1]}；\texttt{jump\_\allowbreak{}\allowbreak{}sel = 11}（\texttt{mepc}） \\
3 & \texttt{dcache\_\allowbreak{}\allowbreak{}stall} & \texttt{stall[4:0]} 全冻结 \\
4 & \texttt{dmem\_\allowbreak{}\allowbreak{}axi\_\allowbreak{}\allowbreak{}stall\_\allowbreak{}\allowbreak{}req}（store 等待总线） & \texttt{stall[4:0]} 全冻结 \\
5 & \texttt{ex\_\allowbreak{}\allowbreak{}stall\_\allowbreak{}\allowbreak{}req}（多周期单元与 ISE） & \texttt{stall[4:2]}；\texttt{flush[1]}（EX/MEM 插气泡） \\
6 & \texttt{jump\_\allowbreak{}\allowbreak{}en}（分支重定向） & \texttt{flush[3:2]}（IF/ID 与 ID/EX）；\texttt{jump\_\allowbreak{}\allowbreak{}sel = 01} \\
7 & \texttt{load\_\allowbreak{}\allowbreak{}use\_\allowbreak{}\allowbreak{}rs1}／\texttt{load\_\allowbreak{}\allowbreak{}use\_\allowbreak{}\allowbreak{}rs2} & \texttt{flush[2]}（ID/EX 插气泡）；\texttt{stall[4:3]}（IF 与 IF/ID） \\
8 & \texttt{icache\_\allowbreak{}\allowbreak{}stall}／\texttt{fence\_\allowbreak{}\allowbreak{}i\_\allowbreak{}\allowbreak{}stall} & \texttt{stall[4:3]}；\texttt{flush[2]} \\
\bottomrule
\end{longtable}
}

表中次序为 \texttt{ctrl.\allowbreak{}v} 内 if-else 链的实际判定顺序，据此，Cache 与总线停顿优先于分支重定向，中断与陷入同属首个条件分支。

该层为加速指令正确执行的关键：ISE 的多拍执行属第 5 项，陷入与中断属第 1 项；两者叠加时以 ``陷入优先、ISE 状态一并清零、不发布迟到结果''为语义（§2.3.3.4）。

\Needspace{85pt}

\subsection{与加速及时序收敛相关的结构改造}

{\small
\renewcommand{\thetable}{2-5}
\begin{longtable}{@{}>{\raggedright\arraybackslash}p{\dimexpr 0.0688\linewidth-2\tabcolsep\relax}>{\raggedright\arraybackslash}p{\dimexpr 0.3480\linewidth-2\tabcolsep\relax}>{\raggedright\arraybackslash}p{\dimexpr 0.2434\linewidth-2\tabcolsep\relax}>{\raggedright\arraybackslash}p{\dimexpr 0.3398\linewidth-2\tabcolsep\relax}@{}}
\caption{与加速及时序收敛相关的结构改造}\label{tab:2-5}\\
\toprule
\textbf{序号} & \textbf{改造内容} & \textbf{位置} & \textbf{动机} \\
\midrule
\endfirsthead
\multicolumn{4}{@{}l}{\footnotesize 续表 2-5}\\
\toprule
\textbf{序号} & \textbf{改造内容} & \textbf{位置} & \textbf{动机} \\
\midrule
\endhead
1 & 器件与工程迁移：Kintex-7 $\rightarrow$ Zynq-7020，重建 3 个 IP 与 1 个 XDC，PLL 由 50 MHz 倍频至 100 MHz & \texttt{board/\allowbreak{}vivado\_\allowbreak{}\allowbreak{}zynq7000\_\allowbreak{}\allowbreak{}xc7z020\_\allowbreak{}\allowbreak{}ascon/\allowbreak{}} & 早期工程的目标器件为 \texttt{xc7k325t}，与提交板卡不符 \\
2 & D-cache 增加 \texttt{MISS\_\allowbreak{}\allowbreak{}CAPTURE} 状态，miss 属性锁存一拍后再发总线 & \texttt{dcache.\allowbreak{}v} & 消除"miss 判定至总线请求"的长组合路径（11 个失败端点） \\
3 & store 与 MMIO 绕过 \texttt{MISS\_\allowbreak{}\allowbreak{}CAPTURE}，恢复单拍注册请求 & \texttt{dcache.\allowbreak{}v} & \texttt{MISS\_\allowbreak{}\allowbreak{}CAPTURE} 对不可 Cache 访问仅增加延迟，据此回收周期代价 \\
4 & 乘法器输入与输出寄存器化，并引入 \texttt{mul\_\allowbreak{}\allowbreak{}ack}、\texttt{cancel} 握手 & \texttt{mult.\allowbreak{}v}、\texttt{ex\_\allowbreak{}\allowbreak{}stage.\allowbreak{}v} & 乘法器组合深度对应 48 个失败端点、$-$0.105 ns \\
5 & ISE 状态寄存器改为局部复位控件 & \texttt{ascon/\allowbreak{}ascon\_\allowbreak{}\allowbreak{}ise.\allowbreak{}v} & 全局复位扇出对应 93 个失败端点、$-$0.128 ns \\
6 & 新增 Ascon MMIO 窗口与 ISE 指令 & \texttt{ascon/\allowbreak{}}、\texttt{rtl\_\allowbreak{}\allowbreak{}p0/\allowbreak{}soc\_\allowbreak{}\allowbreak{}top.\allowbreak{}v} & 当前短消息加速的核心功能 \\
7 & UART 起始位检测审计（\texttt{d0} 与 \texttt{d1/\allowbreak{}d2}） & \texttt{uart\_\allowbreak{}\allowbreak{}rx.\allowbreak{}v} & 已记录同步链风险；未做错误概率实测，不以估计数值判定无影响 \\
\bottomrule
\end{longtable}
}

前五项与加速功能无直接关联，但决定该 SoC 能否在 100 MHz 收敛。第 4 章将说明：33 组策略扫掠将 WNS 自 $-$0.942 ns 提升至 $-$0.118 ns，而剩余 0.118 ns 须由上述结构改造回收。

\Needspace{85pt}

\section{硬件系统介绍}

\Needspace{85pt}

\subsection{硬件整体介绍}

当前入口为 \texttt{board/\allowbreak{}vivado\_\allowbreak{}zynq7000\_\allowbreak{}xc7z020\_\allowbreak{}pynq/\allowbreak{}build.\allowbreak{}tcl}，top=\texttt{pynq\_\allowbreak{}wrapper}、 part=\texttt{xc7z020clg400-\allowbreak{}1}、Vivado 2025.2。PS7 GP0 连接两个 AXI-Lite 窗口， FCLK0 配置为 100 MHz；\texttt{proc\_\allowbreak{}sys\_\allowbreak{}reset} 与控制窗口管理复位。 COE 作为 128 KiB BRAM 初值，由服务固件 ELF/BIN 生成；板上不做动态固件装载。

源集明确使用 \texttt{rtl\_\allowbreak{}p0/\allowbreak{}soc\_\allowbreak{}top.\allowbreak{}v} 和 \texttt{rtl\_\allowbreak{}p0/\allowbreak{}ascon\_\allowbreak{}mmio.\allowbreak{}v}，其余 CPU/Ascon 模块按构建脚本选择。实际逐文件哈希、IP 版本与产物大小见 release/manifest.json。不能用当前 main 的同名文件、旧工程缓存或不同 COE 替代该发布源集。 \texttt{rebuild\_\allowbreak{}firmware.\allowbreak{}tcl} 用于已有独立工程的 BRAM 重建与后续实现，不等于干净构建入口。

历史 PL-only 基准使用 \texttt{top\_\allowbreak{}pl\_\allowbreak{}bram}、板外串口及 clk\_wiz\_0；它与当前 PS 时钟工程分开保存。三套历史固件对应三份 bitstream，当前 PYNQ 则在同一 overlay 下按 mode 选择三方案。

\figpdf{2-6}{fig_2_2}{历史 PL-only 工程结构}{top\_pl\_bram 的板外时钟与串口入口， 仅用于解释历史基准；当前 PS7/FCLK0/AXI-Lite 结构见 PS/PL 通路图。}

\Needspace{85pt}

\subsection{机械设计介绍}

本作品为纯电子系统，不含机械结构与运动部件。板卡采用正点原子领航者 ZYNQ7020 开发板，以底板 + 核心板的两段式结构装配；设计成果全部位于 PL 内部逻辑与 PS 端软件，不需要额外的机械加工件、传感器支架或外部供电模块。板级连接的仅有 50 MHz 时钟源、 USB 串口（UART）、两个用户按键、两个用户 LED 与以太网口，其中串口与网口分别用于固件日志输出与 PYNQ 远程访问。

\Needspace{85pt}

\subsection{电路各模块介绍}

本节介绍三种密码实现和当前 PS/PL 接口；处理器核见 2.2 节。当前源集选择以发布 manifest 与 build.tcl 为准；同名源文件的其他变体不自动属于本次交付。职责、接口、测量范围与验证边界分别说明，不以设计意图替代已取得的证据。

\Needspace{85pt}

\paragraph{2.3.3.1 算法标准与参数定义}

Ascon 的安全核心为 320 bit 状态 \(S=(s_0,s_1,s_2,s_3,s_4)\)，每个分量 64 bit，认证加密通过''吸收与挤出＋置换''实现。AEAD128 的参数如下表。

{\small
\renewcommand{\thetable}{2-6}
\begin{longtable}{@{}>{\raggedright\arraybackslash}p{\dimexpr 0.1895\linewidth-2\tabcolsep\relax}>{\raggedright\arraybackslash}p{\dimexpr 0.4884\linewidth-2\tabcolsep\relax}>{\raggedright\arraybackslash}p{\dimexpr 0.3221\linewidth-2\tabcolsep\relax}@{}}
\caption{Ascon-AEAD128 参数}\label{tab:2-6}\\
\toprule
\textbf{参数} & \textbf{取值} & \textbf{本作品的处理} \\
\midrule
\endfirsthead
\multicolumn{3}{@{}l}{\footnotesize 续表 2-6}\\
\toprule
\textbf{参数} & \textbf{取值} & \textbf{本作品的处理} \\
\midrule
\endhead
key & 16 byte / 128 bit & $K_0$ 取低 64 bit，$K_1$ 取高 64 bit \\
nonce & 16 byte / 128 bit & $N_0$、$N_1$ 同理 \\
tag & 16 byte / 128 bit & 完整输出，不截断 \\
rate / capacity & 128 bit / 192 bit & 每块吸收 16 byte \\
初始向量 & 整数 \texttt{0x0000\_\allowbreak{}\allowbreak{}1000\_\allowbreak{}\allowbreak{}808c\_\allowbreak{}\allowbreak{}0001}，字节串 \texttt{01 00 8c 80 00 10 00 00} & 常量以位移式匹配 \\
轮常量 & \texttt{f0 e1 d2 c3 b4 a5 96 87 78 69 5a 4b}（第 0 至 11 轮） & P12 取前 12 个，P8 取后 8 个 \\
域分离 & $s_4 \mathrel{\oplus}= 2^{63}$（\texttt{0x8000\_\allowbreak{}\allowbreak{}0000\_\allowbreak{}\allowbreak{}0000\_\allowbreak{}\allowbreak{}0000}） & 于 AD 与 PT 之间插入 \\
\bottomrule
\end{longtable}
}

\Needspace{174pt}

\textbf{轮函数}。第 \(i\) 轮先执行常量加与 S-box，再执行线性扩散：

\[
\begin{aligned}
s_2 &\mathrel{\oplus}= c_i \\
s_0 &\mathrel{\oplus}= s_4,\quad s_4 \mathrel{\oplus}= s_3,\quad s_2 \mathrel{\oplus}= s_1 \\
t_0 &= \overline{s_0}\, s_1,\;\; t_1 = \overline{s_1}\, s_2,\;\; t_2 = \overline{s_2}\, s_3,\;\; t_3 = \overline{s_3}\, s_4,\;\; t_4 = \overline{s_4}\, s_0 \\
s_0 &\mathrel{\oplus}= t_1,\quad s_1 \mathrel{\oplus}= t_2,\quad s_2 \mathrel{\oplus}= t_3,\quad s_3 \mathrel{\oplus}= t_4,\quad s_4 \mathrel{\oplus}= t_0 \\
s_1 &\mathrel{\oplus}= s_0,\quad s_0 \mathrel{\oplus}= s_4,\quad s_3 \mathrel{\oplus}= s_2,\quad s_2 = \overline{s_2}
\end{aligned}
\]

\Needspace{93pt}

线性层（\(\mathrm{rotr}\) 表示 64 bit 循环右移）：

\[
s_i \mathrel{\oplus}= \mathrm{rotr}(s_i,\,r_{i,1}) \mathrel{\oplus} \mathrm{rotr}(s_i,\,r_{i,2}),
\qquad (r_{i,1},r_{i,2}) \in \{(19,28),(61,39),(1,6),(10,17),(7,41)\}
\]

\textbf{AEAD 流程}（加密方向）：

\begin{enumerate}
\def\labelenumi{\arabic{enumi}.}
\tightlist
\item
  初始化：\(S \gets (IV,K_0,K_1,N_0,N_1)\)，\(S \gets P^{12}(S)\)，\(s_3 \mathrel{\oplus}= K_0\)， \(s_4 \mathrel{\oplus}= K_1\)；
\item
  关联数据：按 16 byte 分块吸收并执行 \(P^8\)；非空整块之后仍须追加一个空 padding 块并再次执行 \(P^8\)；
\item
  域分离：\(s_4 \mathrel{\oplus}= 2^{63}\)；
\item
  明文：按 16 byte 分块挤出 \(C_i = P_i \oplus (s_0 \| s_1)\) 并执行 \(P^8\)；末块不足 16 byte 时先补 \texttt{0x01} 再补 \texttt{0x00}，末块不再置换；
\item
  终结：\(s_2 \mathrel{\oplus}= K_0\)，\(s_3 \mathrel{\oplus}= K_1\)，\(S \gets P^{12}(S)\)， \(s_3 \mathrel{\oplus}= K_0\)，\(s_4 \mathrel{\oplus}= K_1\)；
\item
  标签：\(T = \mathrm{LE64}(s_3) \,\|\, \mathrm{LE64}(s_4)\)；密文与明文等长。
\end{enumerate}

三种方案遵循的实现要点如下。完整 1,089 组 KAT 为 host 验证范围；当前板级三模式仅验证容量内 289 组：

\begin{itemize}
\tightlist
\item
  padding 为''追加字节 \texttt{0x01} 后补 \texttt{0x00}``，在 64-bit 字中表示为 \(1 \ll (8n)\)，\(n<8\)；
\item
  AD 为空时既不吸收亦不 padding，但仍执行域分离；
\item
  整块 AD 与整块 PT 之后均须再经一个空 padding 块，不得以''满块即结束''处理；
\item
  长度须显式给出，不得由''最高非零 bit''反推（否则 \texttt{00} 与 \texttt{a500} 将判定错误）。
\end{itemize}

\figpdf{2-7}{fig_2_6}{Ascon 置换数据通路}{5 $\times$ 64-bit 状态、S-box 的五个与非项 组合、线性层的五组双旋转与轮常量注入；标注 P12/P8 轮数差异及 ISE 内 320-bit 状态寄存器位置。}

\Needspace{85pt}

\paragraph{2.3.3.2 方案一：纯软件实现（SW）}\label{ux65b9ux6848ux4e00ux7eafux8f6fux4ef6ux5b9eux73b0sw}

软件实现为正确性基准，同时作为周期加速比的参照基准（分母）。生产代码位于 \texttt{ascon/\allowbreak{}ascon\_\allowbreak{}c/\allowbreak{}aead/\allowbreak{}}，锁定且不可修改的参考实现位于 \texttt{sim/\allowbreak{}ascon/\allowbreak{}reference/\allowbreak{}ascon\_\allowbreak{}c/\allowbreak{}}（ascon-c 固定 commit \texttt{446347f2}）。

\textbf{API 契约}（C）：

\begin{itemize}
\tightlist
\item
  \texttt{crypto\_\allowbreak{}aead\_\allowbreak{}encrypt(c,\ clen,\ m,\ mlen,\ ad,\ adlen,\ nsec,\ npub,\ k)}：\texttt{clen} 含 16 byte tag；
\item
  \texttt{crypto\_\allowbreak{}aead\_\allowbreak{}decrypt(m,\ mlen,\ nsec,\ c,\ clen,\ ad,\ adlen,\ npub,\ k)}：tag 拼接于密文末尾，成功返回 0，且 \texttt{*mlen\ =\ clen\ −\ 16}；认证失败返回 $-$1，\texttt{*mlen\ =\ 0}，并清除 \texttt{m{[}0,\ clen−16)} 范围内的候选明文；\texttt{clen\ \textless{}\ 16} 时返回 $-$1 且不写 \texttt{m}。
\end{itemize}

``认证失败须清除候选明文''为项目验收要求，非 SUPERCOP 的强制项：原始 C 参考在 bad tag 时返回 $-$1，但缓冲区仍保留明文，且 \texttt{mlen} 等于明文长度。本作品修复生产 C 实现的失败处理，并将''原始 C 的真实 FAIL 行为''作为保留证据，不予掩盖。

\textbf{历史 PL-only 构建}：以 xpack \texttt{riscv-\allowbreak{}none-\allowbreak{}elf-\allowbreak{}gcc\ 15.\allowbreak{}2.\allowbreak{}0}、\texttt{-\allowbreak{}O2} 交叉编译为裸机固件，链接至 \texttt{0x1000\_\allowbreak{}0000}，经 UART 输出测试结果。周期加速比的分母即该固件在板上的 \texttt{mcycle} 增量。在 16 字节明文、 warm 条件（指令已在 I-cache）下，板上实测为 \textbf{20806 周期}。当前 PYNQ 服务固件的固定输入加密中位周期为 \textbf{23342}，见 §3.3.4.1；不与历史工具链和测量边界混用。

\Needspace{85pt}

\paragraph{2.3.3.3 方案二：存储器映射加速器（MMIO／HW）}\label{ux65b9ux6848ux4e8cux5b58ux50a8ux5668ux6620ux5c04ux52a0ux901fux5668mmiohw}

MMIO 实现将整个 AEAD 实现为挂接于 \texttt{0x3000\_\allowbreak{}0000} 的从设备，CPU 仅承担参数搬运、状态轮询与结果读取。当前接口取自发布源集 \texttt{board/\allowbreak{}vivado\_\allowbreak{}zynq7000\_\allowbreak{}xc7z020\_\allowbreak{}pynq/\allowbreak{}rtl\_\allowbreak{}p0/\allowbreak{}ascon\_\allowbreak{}mmio.\allowbreak{}v}，如下表。

{\small
\renewcommand{\thetable}{2-7}
\begin{longtable}{@{}>{\raggedright\arraybackslash}p{\dimexpr 0.1577\linewidth-2\tabcolsep\relax}>{\raggedright\arraybackslash}p{\dimexpr 0.1577\linewidth-2\tabcolsep\relax}>{\raggedright\arraybackslash}p{\dimexpr 0.0894\linewidth-2\tabcolsep\relax}>{\raggedright\arraybackslash}p{\dimexpr 0.5952\linewidth-2\tabcolsep\relax}@{}}
\caption{MMIO 加速器寄存器}\label{tab:2-7}\\
\toprule
\textbf{偏移} & \textbf{名称} & \textbf{读写} & \textbf{位定义与语义} \\
\midrule
\endfirsthead
\multicolumn{4}{@{}l}{\footnotesize 续表 2-7}\\
\toprule
\textbf{偏移} & \textbf{名称} & \textbf{读写} & \textbf{位定义与语义} \\
\midrule
\endhead
\texttt{0x00} & CTRL & W & \texttt{[0]} START（自清脉冲，仅空闲态接受）；\texttt{[1]} ENCRYPT（1 加密，0 解密）；\texttt{[2]} ZEROIZE（仅空闲态接受，同一写中优先于 START） \\
\texttt{0x04} & STATUS & R & \texttt{[0]} BUSY；\texttt{[1]} DONE；\texttt{[2]} FAIL（解密 tag 校验失败） \\
\texttt{0x08} & ADLEN & W & AD 字节数 \\
\texttt{0x0C} & MLEN & W & 消息字节数 \\
\texttt{0x10-\allowbreak{}-\allowbreak{}0x1C} & KEY & W & 四个小端序 32-bit key 字，由调用者写入 \\
\texttt{0x20-\allowbreak{}-\allowbreak{}0x2C} & NONCE & W & $N_0$ 至 $N_3$ \\
\texttt{0x30-\allowbreak{}-\allowbreak{}0x3C} & TAGIN & W & 解密时待校验的 tag \\
\texttt{0x40} & AD\_\allowbreak{}WDATA & W & 推入一个 AD 字（内部指针自增） \\
\texttt{0x44} & MSG\_\allowbreak{}WDATA & W & 推入一个消息字（加密为明文，解密为密文） \\
\texttt{0x48} & OUT\_\allowbreak{}RDATA & R & 弹出一个输出字（加密为密文，解密为明文），指针自增 \\
\texttt{0x4C-\allowbreak{}-\allowbreak{}0x58} & TAGOUT & R & 加密计算所得 tag \\
\bottomrule
\end{longtable}
}

以下四项行为构成''可用外设''而非''可计算模块''的判据，须写入契约：

\begin{enumerate}
\def\labelenumi{\arabic{enumi}.}
\tightlist
\item
  \textbf{BUSY 期间全部写操作被忽略}，包括 ZEROIZE、START、key/nonce、长度与输入缓冲；
\item
  \textbf{输出读门禁}：\texttt{OUT\_\allowbreak{}RDATA} 仅在 \texttt{DONE\ \&\&\ !BUSY\ \&\&\ !FAIL} 且未超出有效字数时返回数据并推进指针，其余情形返回零且不推进。认证失败不得发布明文由硬件保证，不依赖软件自觉；
\item
  \textbf{ZEROIZE}：清除 key、nonce、tagin、长度与状态，以及全部 AD、MSG、OUT 数组，并经子核心复位清除适配器缓冲、\(s_0..s_4\)、候选明文与 P8/P12 私有副本。DONE 本身不等于清理，取完结果后须显式清零；
\item
  \textbf{长度上限}：硬件实际支持 AD 与消息各 0 至 16 byte（参数名 \texttt{MAX\_\allowbreak{}AD\_\allowbreak{}BYTES=256} 仅为缓冲上限，不代表多块能力），超长请求显式返回 FAIL。该限制为 P0 的边界条件：板级 MMIO 周期加速比仅对 $\leq$16 byte 成立，正文不将其泛化为''任意长度''。
\end{enumerate}

延迟结构上，MMIO 端到端周期由参数搬运（key、nonce、AD 与消息的逐字写入）、轮询往返与密码核本体三部分组成。历史 PL-only 16 B 调用为 \textbf{688 周期}，核本体观测 \textbf{103 周期}；当前 PYNQ 固定输入加密为 \textbf{1439 周期}。核外开销的历史分解不直接套用到当前调用。

\Needspace{85pt}

\paragraph{2.3.3.4 方案三：定制指令扩展（ISE）}\label{ux65b9ux6848ux4e09ux5b9aux5236ux6307ux4ee4ux6269ux5c55ise}

ISE 将置换并入 CPU：320 bit Ascon 状态作为执行级的专用寄存器堆，由 4 条定制指令访问。指令占用 RISC-V 预留的 custom-0 编码（\texttt{0b0001011}），\textbf{未新造主操作码值}，4 条指令以 funct3 区分。

{\small
\renewcommand{\thetable}{2-8}
\begin{longtable}{@{}>{\raggedright\arraybackslash}p{\dimexpr 0.1122\linewidth-2\tabcolsep\relax}>{\raggedright\arraybackslash}p{\dimexpr 0.1491\linewidth-2\tabcolsep\relax}>{\raggedright\arraybackslash}p{\dimexpr 0.3217\linewidth-2\tabcolsep\relax}>{\raggedright\arraybackslash}p{\dimexpr 0.1122\linewidth-2\tabcolsep\relax}>{\raggedright\arraybackslash}p{\dimexpr 0.3048\linewidth-2\tabcolsep\relax}@{}}
\caption{定制指令编码与语义}\label{tab:2-8}\\
\toprule
\textbf{funct3} & \textbf{助记符} & \textbf{操作数} & \textbf{周期} & \textbf{语义} \\
\midrule
\endfirsthead
\multicolumn{5}{@{}l}{\footnotesize 续表 2-8}\\
\toprule
\textbf{funct3} & \textbf{助记符} & \textbf{操作数} & \textbf{周期} & \textbf{语义} \\
\midrule
\endhead
\texttt{000} & \texttt{ascon.\allowbreak{}w} & \texttt{rs1} 为数据，\texttt{rs2[3:0]} 为索引 & 1 & \texttt{state[32*idx +: 32] = rs1} \\
\texttt{001} & \texttt{ascon.\allowbreak{}rd} & \texttt{rs1[3:0]} 为索引，\texttt{rd} 为结果 & 1 & \texttt{rd = state[32*idx +: 32]} \\
\texttt{010} & \texttt{ascon.\allowbreak{}p12} & --- & 多周期 & \texttt{state = P12(state)} \\
\texttt{011} & \texttt{ascon.\allowbreak{}p8} & --- & 多周期 & \texttt{state = P8(state)} \\
\bottomrule
\end{longtable}
}

状态映射固定为 10 个 32-bit 字：\texttt{s0} 低与高、\texttt{s1} 低与高，直至 \texttt{s4}。非法索引（不在 0 至 9 范围）读返回零、写被忽略；该行为属显式容错约定，不产生非法指令异常。保留 opcode 与 funct 编码的全局审计仍在待办之列，正文不声称已覆盖全部指令空间。

实现细节与生命周期如下：

\begin{itemize}
\tightlist
\item
  \texttt{perm\_\allowbreak{}start} 为电平信号（指令位于 EX 且未完成期间保持有效），内部状态机对其打一拍脉冲，完成后锁存结果；据此，流水线 hold 跨越完成不会导致重复执行；
\item
  \textbf{取消}：陷入与中断经 \texttt{cancel} 丢弃在飞操作并清零 ISE 状态，且不发布迟到结果；普通分支冲刷仅清除较年轻指令，不影响已完成的置换；
\item
  \textbf{局部复位寄存器化}：P12/P8 的 1280-bit 私有复位扇出原挂接于全局复位的组合锥（93 个失败端点、 $-$0.128 ns），改为''取消当拍清架构 \texttt{state}，下一拍清私有 \texttt{s\_\allowbreak{}reg/\allowbreak{}s\_\allowbreak{}out}``，接受与 \texttt{p\_\allowbreak{}start} 同拍释放局部复位，不增加正常执行的置换周期（§3.3.1）；
\item
  \textbf{无上下文保存与恢复 ABI}：陷入取消后不可自中间轮续算，调用者须串行独占或整体重算。当前 bare-metal 自测不启用抢占，真实中断、多任务与嵌套调用下的恢复尚未验证；
\item
  \textbf{软件封装}：以 \texttt{asm\ volatile} 手写 intrinsic（不改动 GCC backend），逐字读写状态并调用 P8/P12；软件侧 AEAD 循环与 MMIO 版本共用同一调用顺序与测试向量，故两版密文与 tag 可逐字节比对。
\end{itemize}

历史 PL-only 16 B 加密为 \textbf{2570 周期}、相对同组软件 8.10$\times$；当前 PYNQ 固定输入为 \textbf{3995 周期}、相对当前 SW 5.843$\times$。ISE 的设计动机是用硬件置换替代软件轮函数；不同固件的差值尚未完成分项归因，不把推测写成已测结论。

\figpdf{2-8}{fig_2_7}{三种实现的调用路径对比}{SW（全部于 CPU 内完成）、 HW（CPU 经寄存器搬运至加速器并轮询）、ISE（CPU 直接执行 custom-0，状态位于 EX 级）；标注三者 的数据搬运量与关键延迟分项（搬运、轮询、算法本体）。}

\Needspace{85pt}

\paragraph{2.3.3.5 三种方案的可比性边界}

{\small
\renewcommand{\thetable}{2-9}
\begin{longtable}{@{}>{\raggedright\arraybackslash}p{\dimexpr 0.2059\linewidth-2\tabcolsep\relax}>{\raggedright\arraybackslash}p{\dimexpr 0.2519\linewidth-2\tabcolsep\relax}>{\raggedright\arraybackslash}p{\dimexpr 0.3267\linewidth-2\tabcolsep\relax}>{\raggedright\arraybackslash}p{\dimexpr 0.2155\linewidth-2\tabcolsep\relax}@{}}
\caption{三种方案的可比性边界}\label{tab:2-9}\\
\toprule
\textbf{维度} & \textbf{SW} & \textbf{MMIO（HW）} & \textbf{ISE} \\
\midrule
\endfirsthead
\multicolumn{4}{@{}l}{\footnotesize 续表 2-9}\\
\toprule
\textbf{维度} & \textbf{SW} & \textbf{MMIO（HW）} & \textbf{ISE} \\
\midrule
\endhead
数据通路 & 全部于 CPU 内逐字运算 & CPU $\rightarrow$ MMIO 总线 $\rightarrow$ 专用核 & CPU EX 级内，状态寄存器 \\
单次 AEAD 的 CPU 动作 & 纯计算 & 写 key/\allowbreak{}nonce/\allowbreak{}AD/\allowbreak{}msg/\allowbreak{}len $\rightarrow$ START $\rightarrow$ 轮询 STATUS $\rightarrow$ 读 OUT/TAG & 写状态 $\rightarrow$ P8/P12 $\rightarrow$ 读状态 \\
当前 PYNQ 加密周期（KAT \#529） & 23342 & 1439 & 3995 \\
当前周期加速比（相对 SW） & 1.000$\times$ & 16.221$\times$ & 5.843$\times$ \\
当前服务长度能力 & PT/AD 各 $\leq$16 B & PT/AD 各 $\leq$16 B & PT/AD 各 $\leq$16 B \\
主要成本 & CPU 软件运算 & 寄存器与窗口逻辑、轮询开销 & 指令译码、状态搬运与 EX 级多周期 \\
面积统计 & 三模式共同占用的整机资源；不拆分独立面积 & 同左 & 同左 \\
时序风险 & 无 & MMIO 完成路径（§3.3.1 最差路径） & EX 级 hold 与状态复位扇出 \\
当前使用方式 & 正确性与周期基线 & CPU 搬运参数并主动轮询，短消息加速 & 软件 AEAD 控制流与硬件置换紧耦合 \\
\bottomrule
\end{longtable}
}

上述结果不宜读作''性能强弱排序''，而应理解为两条技术路线的取舍：MMIO 的周期加速比更高，但其能力受 16 byte 寄存器窗口与轮询开销约束，且 CPU 在轮询期间不可用；ISE 的周期加速比相对较低，然其状态位于流水线内，无 MMIO 总线往返，但当前服务仍统一受 16 B 长度限制。三者置于同一 SoC、同一套向量与同一块板上测量，方能为该取舍提供实测依据。

\evid{证据 docs/\allowbreak{}ascon\_\allowbreak{}key\_\allowbreak{}management.\allowbreak{}md；docs/\allowbreak{}ascon\_\allowbreak{}standard\_\allowbreak{}baseline.\allowbreak{}md；board/\allowbreak{}vivado\_\allowbreak{}zynq7000\_\allowbreak{}xc7z020\_\allowbreak{}pynq/\allowbreak{}rtl\_\allowbreak{}p0/\allowbreak{}ascon\_\allowbreak{}mmio.\allowbreak{}v；src/\allowbreak{}rtl/\allowbreak{}ascon/\allowbreak{}ascon\_\allowbreak{}ise.\allowbreak{}v；data/\allowbreak{}pynq/\allowbreak{}notebook\_\allowbreak{}20261029\_\allowbreak{}130920/\allowbreak{}pynq\_\allowbreak{}warm\_\allowbreak{}benchmark.\allowbreak{}json}

\Needspace{85pt}

\paragraph{2.3.3.6 PS/PL 数据通道方案论证}\label{pspl-ux6570ux636eux901aux9053ux65b9ux6848ux8bbaux8bc1}

任务书初稿拟定的方案为''AXI-Lite 控制寄存器 + 共享 BRAM''。查证冻结 RTL 后，该方案在不修改 \texttt{src/\allowbreak{}rtl} 的前提下不可成立，依据如下：

{\small
\renewcommand{\thetable}{2-10}
\begin{longtable}{@{}>{\raggedright\arraybackslash}p{\dimexpr 0.3686\linewidth-2\tabcolsep\relax}>{\raggedright\arraybackslash}p{\dimexpr 0.2628\linewidth-2\tabcolsep\relax}>{\raggedright\arraybackslash}p{\dimexpr 0.3686\linewidth-2\tabcolsep\relax}@{}}
\caption{PS/PL 数据通道选择依据}\label{tab:2-10}\\
\toprule
\textbf{证据} & \textbf{位置} & \textbf{后果} \\
\midrule
\endfirsthead
\multicolumn{3}{@{}l}{\footnotesize 续表 2-10}\\
\toprule
\textbf{证据} & \textbf{位置} & \textbf{后果} \\
\midrule
\endhead
外设访问全部在 \texttt{bus\_\allowbreak{}\allowbreak{}interconnect} 内部译码（clint、uart、gpio、cpi、ascon），无通向 PL 边界的 MMIO 出口 & \texttt{src/\allowbreak{}rtl/\allowbreak{}bus\_\allowbreak{}\allowbreak{}interconnect.\allowbreak{}v} & 若为 CPU 增加 PS 可见寄存器窗口，必须修改该文件 \\
仅 \texttt{0x1000\_\allowbreak{}\allowbreak{}0000} 至 \texttt{0x100F\_\allowbreak{}\allowbreak{}FFFF} 被判为可 Cache & \texttt{src/\allowbreak{}rtl/\allowbreak{}dcache.\allowbreak{}v} & 共享 BRAM 必落于该窗口内，PS 写入后 CPU 可能读到旧 Cache 行 \\
无 D-cache 失效或清洗指令（仅 \texttt{fence.\allowbreak{}i} 触发 I-cache 冲刷） & \texttt{src/\allowbreak{}rtl/\allowbreak{}cpu\_\allowbreak{}\allowbreak{}core.\allowbreak{}v}、\texttt{dcache.\allowbreak{}v} & 固件无法在两次事务之间保证一致性，\texttt{volatile} 与 \texttt{fence} 不足以替代 \\
历史 P0 时序余量为 WNS +0.087 ns & \texttt{docs/\allowbreak{}freeze\_\allowbreak{}\allowbreak{}checkpoint\_\allowbreak{}\allowbreak{}20261027.\allowbreak{}md} §4 & 改动公共路径需重新验证；采用独立 PYNQ 工程控制变更范围 \\
\bottomrule
\end{longtable}
}

据此，本作品采用''以既有 UART MMIO 作为 PS/PL 数据通道''的方案：在 PL 内新增 \texttt{pynq\_\allowbreak{}mbox.\allowbreak{}v} 实现串口的另一端，字节经 \texttt{uart\_\allowbreak{}tx} 送入 SoC 的 \texttt{uart\_\allowbreak{}rxd}，SoC 的 \texttt{uart\_\allowbreak{}txd} 经 \texttt{uart\_\allowbreak{}rx} 收入 RX FIFO 供 PS 读取。该方案具备四项性质：

\begin{enumerate}
\def\labelenumi{\arabic{enumi}.}
\tightlist
\item
  \textbf{独立发布源集}：历史 PL-only 证据保留，当前 BD 与桥接 RTL 单独构建并验证；
\item
  UART 属不可 Cache MMIO，当前无 PS/RV32 共享可缓存 buffer，缓存可见性记为 N/A，不泛化为所有架构均无一致性问题；
\item
  同一份服务固件既可运行于 P0 纯 PL 版本（接真实 UART 引脚），亦可运行于 PYNQ 版本（PS 为对端），仅需维护一条代码路径；
\item
  代价为延迟与带宽：链路为 6.25 Mbaud（\texttt{BAUD\_\allowbreak{}CNT\_\allowbreak{}MAX\ =\ 16}，100 MHz 分频），每字节约 1.6 $\mu$s。理论字节时间不等于 PS 调用墙钟。固定 KAT \#529 的实测中位为 \textbf{4.262--5.216 ms}，包含 Python、MMIO 与 UART。固件 AEAD 周期不含串口，须与墙钟分列；不同固件边界仍会影响周期。
\end{enumerate}

\Needspace{170pt}

\paragraph{2.3.3.7 结构与寄存器定义}

\begin{codeblock}[samepage=true]
Jupyter (Python)
  └─ AsconHost → MailboxBackend → PynqUartTransport（src/pynq/pynq_transport.py）
        │  pynq.MMIO × 2（GP0，各 4 KiB）
        ├─ CONTROL 0x4000_0000 ── pynq_ctl.v  ：SOC_RESET / BUTTON / STAT(LED, RST) / ID
        └─ MAILBOX 0x4000_1000 ── pynq_mbox.v ：TX FIFO（64）─uart_tx─┐
                                                                     │ PL 内部，不经引脚
                                        RX FIFO（256）─uart_rx──┐    │
                                                               ▼    ▼
                                            SoC uart_rxd / uart_txd（top 端口，P0 同源集）
                                              cpu_core + dcache + uart_ctrl（0x2000_0000）
                                                        └─ Ascon：SW／MMIO／ISE
\end{codeblock}

\textbf{CONTROL 窗口}（\texttt{pynq\_\allowbreak{}ctl.\allowbreak{}v}，AXI-Lite 32 bit，遵循 WSTRB{[}0{]}）：

{\small
\renewcommand{\thetable}{2-11}
\begin{longtable}{@{}>{\raggedright\arraybackslash}p{\dimexpr 0.0987\linewidth-2\tabcolsep\relax}>{\raggedright\arraybackslash}p{\dimexpr 0.1460\linewidth-2\tabcolsep\relax}>{\raggedright\arraybackslash}p{\dimexpr 0.0987\linewidth-2\tabcolsep\relax}>{\raggedright\arraybackslash}p{\dimexpr 0.6567\linewidth-2\tabcolsep\relax}@{}}
\caption{CONTROL 窗口寄存器}\label{tab:2-11}\\
\toprule
\textbf{偏移} & \textbf{名称} & \textbf{读写} & \textbf{位定义} \\
\midrule
\endfirsthead
\multicolumn{4}{@{}l}{\footnotesize 续表 2-11}\\
\toprule
\textbf{偏移} & \textbf{名称} & \textbf{读写} & \textbf{位定义} \\
\midrule
\endhead
\texttt{0x00} & CTRL & RW & \texttt{[0]} SOC\_\allowbreak{}RESET（1 保持 RV32 复位，上电默认 1）；\texttt{[3:2]} BUTTON[1:0] 透传至 SoC GPIO \\
\texttt{0x04} & STAT & RO & \texttt{[0]} SOC\_\allowbreak{}RST\_\allowbreak{}N（送给 SoC 的复位，1 为运行）；\texttt{[3:2]} LED[1:0] 回读 \\
\texttt{0x08} & ID & RO & \texttt{0x514E\_\allowbreak{}\allowbreak{}5950}（"PYNQ"，小端） \\
\texttt{0x0C} & VERSION & RO & \texttt{0x0000\_\allowbreak{}\allowbreak{}0001} \\
\bottomrule
\end{longtable}
}

复位约定为：\texttt{CTRL{[}0{]}} 上电置 1（SoC 保持复位），PS 确认链路可用后写 0 释放；释放信号经三级同步器进入 SoC 时钟域。

\textbf{MAILBOX 窗口}（\texttt{pynq\_\allowbreak{}mbox.\allowbreak{}v}，TX 深度 64、RX 深度 256）：

{\small
\renewcommand{\thetable}{2-12}
\begin{longtable}{@{}>{\raggedright\arraybackslash}p{\dimexpr 0.0987\linewidth-2\tabcolsep\relax}>{\raggedright\arraybackslash}p{\dimexpr 0.1460\linewidth-2\tabcolsep\relax}>{\raggedright\arraybackslash}p{\dimexpr 0.0987\linewidth-2\tabcolsep\relax}>{\raggedright\arraybackslash}p{\dimexpr 0.6567\linewidth-2\tabcolsep\relax}@{}}
\caption{MAILBOX 窗口寄存器}\label{tab:2-12}\\
\toprule
\textbf{偏移} & \textbf{名称} & \textbf{读写} & \textbf{含义} \\
\midrule
\endfirsthead
\multicolumn{4}{@{}l}{\footnotesize 续表 2-12}\\
\toprule
\textbf{偏移} & \textbf{名称} & \textbf{读写} & \textbf{含义} \\
\midrule
\endhead
\texttt{0x00} & TXDATA & W & \texttt{[7:0]} 压入一字节；FIFO 满时延迟写响应 \\
\texttt{0x04} & RXDATA & R & \texttt{[7:0]} 弹出一字节；空时读 0 且不弹出 \\
\texttt{0x08} & TXSTAT & R & \texttt{[15:0]} 深度；\texttt{[16]} 满；\texttt{[17]} 空 \\
\texttt{0x0C} & RXSTAT & R & \texttt{[15:0]} 深度；\texttt{[16]} 满；\texttt{[17]} 空；\texttt{[18]} 溢出（粘滞） \\
\texttt{0x10} & RXCLR & W & \texttt{[0]} 清溢出标志 \\
\texttt{0x14} & TXFLUSH & W & \texttt{[0]} 清空 TX FIFO \\
\texttt{0x18} & CNT & R & \texttt{[15:0]} PS 已压入字节数；\texttt{[31:16]} 已接收字节数（16-bit 回绕） \\
\texttt{0x1C} & ID & R & \texttt{0x4D42\_\allowbreak{}\allowbreak{}4F58}。按 32-bit 小端读出字节序为 \texttt{58 4F 42 4D}（"XOBM"），按大端字节序方为 ASCII "MBOX"；与 CONTROL 窗口 \texttt{0x514E\_\allowbreak{}\allowbreak{}5950}（小端读出 "PYNQ"）的编码方向相反，正文按 RTL 现状如实记录 \\
\bottomrule
\end{longtable}
}

\texttt{link\_\allowbreak{}flush}（由 CONTROL 的 SOC\_RESET 驱动）为高时清空两侧 FIFO 与计数，保证每次 SoC 复位后链路回到确定状态。

\Needspace{85pt}

\paragraph{2.3.3.8 服务协议与服务固件}

字节流格式由 \texttt{src/\allowbreak{}pynq/\allowbreak{}mailbox\_\allowbreak{}protocol.\allowbreak{}py} 与服务固件共同实现，共四类帧：

{\small
\renewcommand{\thetable}{2-13}
\begin{longtable}{@{}>{\raggedright\arraybackslash}p{\dimexpr 0.0700\linewidth-2\tabcolsep\relax}>{\raggedright\arraybackslash}p{\dimexpr 0.3657\linewidth-2\tabcolsep\relax}>{\raggedright\arraybackslash}p{\dimexpr 0.1026\linewidth-2\tabcolsep\relax}>{\raggedright\arraybackslash}p{\dimexpr 0.4617\linewidth-2\tabcolsep\relax}@{}}
\caption{服务帧类型}\label{tab:2-13}\\
\toprule
\textbf{帧} & \textbf{长度} & \textbf{方向} & \textbf{用途} \\
\midrule
\endfirsthead
\multicolumn{4}{@{}l}{\footnotesize 续表 2-13}\\
\toprule
\textbf{帧} & \textbf{长度} & \textbf{方向} & \textbf{用途} \\
\midrule
\endhead
\texttt{Y01Q} & 48 B 头 + key16 + nonce16 + AD + data + tag & PS $\rightarrow$ PL & AEAD 请求 \\
\texttt{Y01R} & 48 B 头 + payload & PL $\rightarrow$ PS & AEAD 响应 \\
\texttt{Y01S} & 64 B & PL $\rightarrow$ PS & 开机与 epoch 应答：固件标识、能力上限、自检结果、boot \texttt{mcycle} \\
\texttt{Y01E} & 20 B & PS $\rightarrow$ PL & PS 下发会话 epoch \\
\bottomrule
\end{longtable}
}

请求头字段依次为：magic（\texttt{Y01Q}）、version、flags、epoch（u64）、request\_id（u64）、operation （1 加密 / 2 解密）、mode（1 SW / 2 HW / 3 ISE）、adlen、dlen、output\_capacity、taglen；响应头为 magic（\texttt{Y01R}）、version、flags、epoch、request\_id、status、operation、mode、size、cycles。状态码为 \texttt{OK}、\texttt{AUTH\_\allowbreak{}FAILED}、\texttt{LENGTH}、\texttt{UNSUPPORTED}、\texttt{INTERNAL}；响应 flags 中 \texttt{SIMULATED} 位在真实硬件执行时恒为 0，\texttt{CYCLES\_\allowbreak{}VALID} 位指示周期字段是否有效。协议约束为：失败响应不得携带数据与周期；key 与 nonce 长度须为 16；解密请求须携带 16 字节 tag。

服务固件 \texttt{pynq-\allowbreak{}svc-\allowbreak{}v1}（\texttt{src/\allowbreak{}firmware/\allowbreak{}pynq\_\allowbreak{}svc/\allowbreak{}}）具备以下设计特征：

\begin{itemize}
\tightlist
\item
  \textbf{单槽串行}：一次仅处理一个请求，响应发完后回到帧同步状态，故 PS 侧无需额外 ACK 线；
\item
  \textbf{周期口径}：仅在 AEAD 调用边界读取 64-bit \texttt{mcycle}（\texttt{0xB00}/\texttt{0xB80}，hi-lo-hi 一致读取），\textbf{不含串口传输时间}；
\item
  \textbf{能力上限}：message 与 AD 各 16 字节，超长请求立即拒绝并返回 \texttt{LENGTH}；
\item
  \textbf{构建}：\texttt{src/\allowbreak{}firmware/\allowbreak{}pynq\_\allowbreak{}svc/\allowbreak{}build.\allowbreak{}sh} 以 RV32IM\_Zicsr/ilp32 交叉编译，产物中检出 24 条 custom-0 指令，P12/P8 均存在，可作为''ISE 确实被调用''的静态证据；
\item
  \textbf{异常处置}：半帧读取依赖主机超时后 reset；无本地接收超时；UART 无远端流控。
\end{itemize}

\figpdf{2-9}{fig_2_9}{PS/PL 协同通路结构}{Jupyter $\rightarrow$ Python API $\rightarrow$ 两个 AXI-Lite 窗口 $\rightarrow$ PL 内 UART 桥 $\rightarrow$ RV32 服务固件 $\rightarrow$ SW/MMIO/ISE；标注时钟域（FCLK0 100 MHz）、FIFO 深度与 帧格式。}

\Needspace{85pt}

\paragraph{2.3.3.9 Overlay 构建与时序签核}\label{overlay-ux6784ux5efaux4e0eux65f6ux5e8fux7b7eux6838}

Overlay 由 \texttt{board/\allowbreak{}vivado\_\allowbreak{}zynq7000\_\allowbreak{}xc7z020\_\allowbreak{}pynq/\allowbreak{}build.\allowbreak{}tcl} 自动创建：独立工程 \texttt{pynq\_\allowbreak{}overlay.\allowbreak{}xpr}、 Block Design（PS7 GP0、FCLK0 = 100 MHz、\texttt{proc\_\allowbreak{}sys\_\allowbreak{}reset}）、综合与实现并导出同名 bit 与 hwh。发布产物为 \texttt{release/\allowbreak{}pynq.\allowbreak{}bit}（4,045,688 B，与 P0 三份 4,045,687 B 的 bitstream 属不同源集）与 \texttt{release/\allowbreak{}pynq.\allowbreak{}hwh}（148,833 B），二者由同一 BD 导出；BRAM 单独重新生成后，MIF 与最终 BIN 逐字节一致（128 KiB，含补零）。

\figpdf{2-10}{fig_2_13}{Vivado Block Design：PYNQ overlay}{在 `pynq\_overlay` 工程中 打开 Block Design，展示 PS7、AXI Interconnect、`proc\_sys\_reset` 与 PL 内 CONTROL / MAILBOX 窗口及 SoC 的连接。}

时序与实现结果（100 MHz，2026-10-28）：

{\small
\renewcommand{\thetable}{2-14}
\begin{longtable}{@{}>{\raggedright\arraybackslash}p{\dimexpr 0.1248\linewidth-2\tabcolsep\relax}>{\raggedright\arraybackslash}p{\dimexpr 0.4376\linewidth-2\tabcolsep\relax}>{\raggedright\arraybackslash}p{\dimexpr 0.4376\linewidth-2\tabcolsep\relax}@{}}
\caption{Overlay 时序与实现结果}\label{tab:2-14}\\
\toprule
\textbf{指标} & \textbf{数值} & \textbf{判据} \\
\midrule
\endfirsthead
\multicolumn{3}{@{}l}{\footnotesize 续表 2-14}\\
\toprule
\textbf{指标} & \textbf{数值} & \textbf{判据} \\
\midrule
\endhead
WNS / TNS & \textbf{+0.034 ns / 0} & setup 失败端点 0/62861 \\
WHS / THS & +0.020 ns / 0 & hold 失败端点 0/62861 \\
未约束检查 & 内部未约束端点 0；\texttt{no\_\allowbreak{}\allowbreak{}clock}、\texttt{no\_\allowbreak{}\allowbreak{}input\_\allowbreak{}\allowbreak{}delay}、\texttt{no\_\allowbreak{}\allowbreak{}output\_\allowbreak{}\allowbreak{}delay} 均为 0 & 本工程无外接 PL IO，故未套用 PL-only 的 UART/LED XDC \\
DRC & \textbf{0 Error / 0 Critical Warning} & 默认显示 53 项 Warning，含规则条数上限提示，不等于全部潜在 Warning 数 \\
整机资源 & LUT 14269、FF 19691、BRAM 34.5 tile、DSP 4 & utilization.rpt；不沿用历史 PL-only 数字 \\
\bottomrule
\end{longtable}
}

\figpdf{2-11}{fig_2_14}{Vivado 实现后时序摘要}{打开 `impl\_1` 的 Report Timing Summary， 截取 Design Timing Summary 区域，应可见 WNS +0.034 ns、TNS 0、WHS +0.020 ns 与失败端点 0。}

须注意：该 overlay 工程使用 \texttt{rtl\_\allowbreak{}p0} 目录下的 \texttt{soc\_\allowbreak{}top.\allowbreak{}v} 与 \texttt{ascon\_\allowbreak{}mmio.\allowbreak{}v} 副本，其余 CPU 与 Ascon 模块按 manifest 记录选择。当前 KEY 由调用者写入，不能把工作树任意同名文件当作发布输入。本轮不改 RTL、固件、bit/hwh 或 manifest，只对齐报告与已有证据。

\Needspace{85pt}

\paragraph{2.3.3.10 验证结果与板上实测}

\textbf{仿真与主机侧验证}（\texttt{validation\_\allowbreak{}20261028.\allowbreak{}md}，全部 exit 0）：

{\small
\renewcommand{\thetable}{2-15}
\begin{longtable}{@{}>{\raggedright\arraybackslash}p{\dimexpr 0.3011\linewidth-2\tabcolsep\relax}>{\raggedright\arraybackslash}p{\dimexpr 0.6077\linewidth-2\tabcolsep\relax}>{\raggedright\arraybackslash}p{\dimexpr 0.0913\linewidth-2\tabcolsep\relax}@{}}
\caption{仿真与主机侧验证}\label{tab:2-15}\\
\toprule
\textbf{验证项} & \textbf{内容} & \textbf{结果} \\
\midrule
\endfirsthead
\multicolumn{3}{@{}l}{\footnotesize 续表 2-15}\\
\toprule
\textbf{验证项} & \textbf{内容} & \textbf{结果} \\
\midrule
\endhead
\texttt{test\_\allowbreak{}\allowbreak{}host\_\allowbreak{}\allowbreak{}api.\allowbreak{}py} & 17 项，含历史 Notebook 顺序执行 & PASS \\
\texttt{test\_\allowbreak{}\allowbreak{}mailbox.\allowbreak{}py} & 24 项协议编解码与错误分支 & PASS \\
\texttt{test\_\allowbreak{}\allowbreak{}navigator.\allowbreak{}py} & 8 项 & PASS \\
\texttt{test\_\allowbreak{}\allowbreak{}pynq\_\allowbreak{}\allowbreak{}transport.\allowbreak{}py} & 14 项，含 real-service 标志回归 & PASS \\
\texttt{golden\_\allowbreak{}\allowbreak{}checker.\allowbreak{}py} & 1,089 组官方 KAT、4,290 次拒绝（host 范围） & PASS \\
\texttt{run\_\allowbreak{}\allowbreak{}bridge.\allowbreak{}py} & 真实桥 RTL：双向传输、空读、WSTRB、RX 满与溢出、顺序读回、复位中止发送 & PASS \\
\texttt{run\_\allowbreak{}\allowbreak{}service.\allowbreak{}py} & 真实 CPU + 最终固件 + AXI-Lite/UART + 行为级 AXI 存储；三模式 36 请求（加密/解密/坏 tag），检查 \texttt{mcycle} 非零；半帧复位后重新握手 & PASS \\
固件构建 & RV32IM\_\allowbreak{}Zicsr/ilp32，24 条 custom-0，P12/P8 存在 & PASS \\
镜像一致性 & Vivado BRAM MIF 与最终 BIN 逐字节比较（128 KiB 含补零） & PASS \\
\bottomrule
\end{longtable}
}

\textbf{板上结果}（正点原子领航者 ZYNQ7020，PYNQ 2.7.0，Python 3.8.2）： overlay 加载成功，MMIO ID 检查通过；服务报告为 \texttt{firmware\_\allowbreak{}id=pynq-\allowbreak{}svc-\allowbreak{}v1}、\texttt{simulated=false}、 \texttt{selftest\_\allowbreak{}mask=7}、\texttt{max\_\allowbreak{}message\_\allowbreak{}bytes=16}、\texttt{max\_\allowbreak{}ad\_\allowbreak{}bytes=16}、\texttt{board\_\allowbreak{}validation=PASS}。

板上在 16 字节容量内执行 message 与 AD 各 0 至 16 字节的全部 289 组组合，三模式均通过：

{\footnotesize
\renewcommand{\thetable}{2-16}
\begin{longtable}{@{}>{\raggedright\arraybackslash}p{\dimexpr 0.0818\linewidth-2\tabcolsep\relax}>{\raggedright\arraybackslash}p{\dimexpr 0.1328\linewidth-2\tabcolsep\relax}>{\raggedright\arraybackslash}p{\dimexpr 0.1866\linewidth-2\tabcolsep\relax}>{\raggedright\arraybackslash}p{\dimexpr 0.1553\linewidth-2\tabcolsep\relax}>{\raggedright\arraybackslash}p{\dimexpr 0.1328\linewidth-2\tabcolsep\relax}>{\raggedright\arraybackslash}p{\dimexpr 0.1553\linewidth-2\tabcolsep\relax}>{\raggedright\arraybackslash}p{\dimexpr 0.1553\linewidth-2\tabcolsep\relax}@{}}
\caption{板上容量内 KAT 与周期统计}\label{tab:2-16}\\
\toprule
\textbf{模式} & \textbf{加密 KAT} & \textbf{错误 Tag 拒绝} & \textbf{周期最小值} & \textbf{周期均值} & \textbf{周期中位数} & \textbf{周期最大值} \\
\midrule
\endfirsthead
\multicolumn{7}{@{}l}{\footnotesize 续表 2-16}\\
\toprule
\textbf{模式} & \textbf{加密 KAT} & \textbf{错误 Tag 拒绝} & \textbf{周期最小值} & \textbf{周期均值} & \textbf{周期中位数} & \textbf{周期最大值} \\
\midrule
\endhead
SW & 289/289 & 289/289 & 17584 & 23801.7 & 23543 & 33595 \\
HW & 289/289 & 289/289 & 1229 & 1589.0 & 1594 & 1797 \\
ISE & 289/289 & 289/289 & 3012 & 4122.0 & 4137 & 5209 \\
\bottomrule
\end{longtable}
}

据同一统计口径，均值之比为 \textbf{SW/HW $\approx$ 15.0$\times$、SW/ISE $\approx$ 5.8$\times$}，不是逐条加速比的均值。该组数据不能与 §3.3.4.3 的历史 PL-only 独立基准数值直接比较：后者为固定 16 字节、单一拍数的端到端测量（SW 20806、ISE 2570、MMIO 688），而本表为服务固件覆盖 0 至 16 字节全长度组合的逐次统计，且调用内部包含适配、轮询与缓冲清理，不含请求帧传输。两组数值差异的归因（尤以 HW 路径的平均 1589 周期与 P0 的 688 周期之差别）本轮\blockedflag{尚未完成}，故正文分列呈现，不作合并。

\textbf{证据身份}：板上测试使用发布版 bit、hwh 与固件，对应 Git \texttt{ac9143d5}，bit SHA-256 \texttt{6e7f9320…}，hwh SHA-256 \texttt{3ebd2640…}；证据归档于 \texttt{data/\allowbreak{}pynq/\allowbreak{}board\_\allowbreak{}20261028/\allowbreak{}}（\texttt{pynq\_\allowbreak{}board\_\allowbreak{}full.\allowbreak{}json}、 \texttt{pynq\_\allowbreak{}board\_\allowbreak{}full.\allowbreak{}log}、\texttt{pynq\_\allowbreak{}board\_\allowbreak{}20261028.\allowbreak{}log}、\texttt{pynq\_\allowbreak{}board\_\allowbreak{}result.\allowbreak{}json}、\texttt{identity.\allowbreak{}txt}）。

\textbf{10-29 固定输入 benchmark 与 Notebook}：同 release、官方 KAT \#529、PT=16 B、AD=0，每模式预热 5 次后各 30 次加密/解密，全部 KAT 匹配。周期和墙钟见 §3.3.4.1--§3.3.4.2，原始数组见 \texttt{data/\allowbreak{}pynq/\allowbreak{}notebook\_\allowbreak{}20261029\_\allowbreak{}130920/\allowbreak{}}。顺序执行结果存在；warm benchmark 已执行 ipynb 本体仍缺。

\textbf{10-29 V03 错误恢复}：9/9 checks、7/7 reset 恢复通过；epoch 1$\rightarrow$8，每次恢复后重新验证三模式 KAT。17 B PT/AD 返回 LENGTH；半帧、过期 deadline、重复 ID、错误 epoch、 RX FIFO overflow 均按门禁验证；每模式 50 次加密与 50 次解密匹配 KAT。过期 deadline 用例验证主机截止时间与恢复机制，不证明硬件永久卡死也可自动恢复。 \texttt{data/\allowbreak{}pynq/\allowbreak{}error\_\allowbreak{}gate\_\allowbreak{}20261029\_\allowbreak{}134023/\allowbreak{}} 含执行 Notebook、identity、JSON 与日志。

\figpdf{2-12}{fig_2_10}{PYNQ 板测数据重排}{容量内 KAT、固定输入 warm benchmark 与错误恢复， 取自归档 JSON，不是原始终端截图或真实 Python 调用代码示例。}

\Needspace{85pt}

\paragraph{2.3.3.11 未闭环项}

{\small
\renewcommand{\thetable}{2-17}
\begin{longtable}{@{}>{\raggedright\arraybackslash}p{\dimexpr 0.3012\linewidth-2\tabcolsep\relax}>{\raggedright\arraybackslash}p{\dimexpr 0.2410\linewidth-2\tabcolsep\relax}>{\raggedright\arraybackslash}p{\dimexpr 0.4578\linewidth-2\tabcolsep\relax}@{}}
\caption{PYNQ 通路未闭环项}\label{tab:2-17}\\
\toprule
\textbf{项} & \textbf{状态} & \textbf{说明} \\
\midrule
\endfirsthead
\multicolumn{3}{@{}l}{\footnotesize 续表 2-17}\\
\toprule
\textbf{项} & \textbf{状态} & \textbf{说明} \\
\midrule
\endhead
超出容量的 800 组板上用例 & NOT\_\allowbreak{}RUN / 当前接口不支持 & 全量 host 已通过，不列为本次新增功能承诺 \\
warm benchmark 执行 Notebook 本体 & 待归档 & JSON/log 已存在，error-gate 本体已归档 \\
小时级与资源泄漏 & not\_\allowbreak{}measured & 每模式 50 次往返不等于小时级 RSS/FD 统计 \\
协议错误与恢复 & 已有 9/9 实板证据 & 当前无共享 cacheable buffer，不伪造缓存测试 \\
重复断电上电稳定性 & not\_\allowbreak{}measured & 需冷启动与每次 smoke，软件 reset 不等于断电 \\
独立 FCLK 测量 & not\_\allowbreak{}measured & 配置 100 MHz 不等于独立测量 \\
硬件轮询超时的错误语义 & \todoflag{待完善} & 当前可能表现为内部失败或解密认证失败；异常时整链 reset 后重测，不作为认证性能数据 \\
\bottomrule
\end{longtable}
}

\evid{证据 board/\allowbreak{}vivado\_\allowbreak{}zynq7000\_\allowbreak{}xc7z020\_\allowbreak{}pynq/\allowbreak{}validation\_\allowbreak{}20261028.\allowbreak{}md；board/\allowbreak{}vivado\_\allowbreak{}zynq7000\_\allowbreak{}xc7z020\_\allowbreak{}pynq/\allowbreak{}validation\_\allowbreak{}20261029.\allowbreak{}md；data/\allowbreak{}pynq/\allowbreak{}board\_\allowbreak{}20261028/\allowbreak{}pynq\_\allowbreak{}board\_\allowbreak{}full.\allowbreak{}json；data/\allowbreak{}pynq/\allowbreak{}notebook\_\allowbreak{}20261029\_\allowbreak{}130920/\allowbreak{}；data/\allowbreak{}pynq/\allowbreak{}error\_\allowbreak{}gate\_\allowbreak{}20261029\_\allowbreak{}134023/\allowbreak{}}

\Needspace{85pt}

\subsection{独立扩展原型：RO-PUF 与器件绑定通路}\label{ux72ecux7acbux6269ux5c55ux539fux578bro-puf-ux4e0eux5668ux4ef6ux7ed1ux5b9aux901aux8def}

本章保留已实现的进阶探索，范围为独立 PL-only 扩展原型，不属于当前 PYNQ release。其成果是 RO 采样、注册、重建判定、派生和内部接线的功能链，不宣称已经完成产品级硬件信任根或安全启动。

\Needspace{85pt}

\paragraph{2.3.4.1 架构与已实现代码}

\texttt{src/\allowbreak{}rtl/\allowbreak{}keys/\allowbreak{}} 包含 device\_dna、ro\_unit、ro\_puf\_core、puf\_key 和 key\_mmio。 256 路 RO 经计数、两两比较及多次表决产生 128-bit 响应；注册保存参考响应与掩码，重建将当前响应的掩码位与参考比较。Device DNA 提供 57-bit 器件标识，不是秘密熵源。

派生通路将固定上下文、响应及 DNA 组成 320-bit 状态，执行 P12 后取低 128 bit 作为 key。这是项目自定义的置换派生原型，不是 NIST SP 800-232 对该 KDF 构造的认证或安全证明。 key 寄存器本身没有 CPU 读地址，内部连接至扩展版 ascon\_mmio；当前 PYNQ 的调用者 key 路径与此不同。采样窗口由构建参数决定，不使用旧稿的 3.3 ms 估算作为板级实测延迟。

\figpdf{2-13}{fig_2_8}{独立扩展原型结构}{RO 阵列、计数与表决、参考和掩码、 DNA 与 P12 派生、128-bit 内部 key；调试回读通路和未验收时钟边界均保留说明。}

\Needspace{85pt}

\paragraph{2.3.4.2 实板观测与失败保留}

证据为 \texttt{data/\allowbreak{}ascon/\allowbreak{}results/\allowbreak{}board/\allowbreak{}uart\_\allowbreak{}20261027\_\allowbreak{}keys3.\allowbreak{}txt} 及对应 XDat。同一文件含两个版本片段，必须区分：

{\small
\renewcommand{\thetable}{2-18}
\begin{longtable}{@{}>{\raggedright\arraybackslash}p{\dimexpr 0.1629\linewidth-2\tabcolsep\relax}>{\raggedright\arraybackslash}p{\dimexpr 0.4512\linewidth-2\tabcolsep\relax}>{\raggedright\arraybackslash}p{\dimexpr 0.3859\linewidth-2\tabcolsep\relax}@{}}
\caption{实板观测与失败保留}\label{tab:2-18}\\
\toprule
\textbf{片段} & \textbf{原始状态} & \textbf{可以支持的结论} \\
\midrule
\endfirsthead
\multicolumn{3}{@{}l}{\footnotesize 续表 2-18}\\
\toprule
\textbf{片段} & \textbf{原始状态} & \textbf{可以支持的结论} \\
\midrule
\endhead
v2 抢跑片段 & busy 时启动，出现 quality=0，后续 127，最终 FAIL & 不能把该 127/128 当作正式通过结果 \\
v3 先等 idle & enroll=0x00008026，quality=128；MASK 全 1、DITHER 全 0、REF=RAW & 此次注册 128/128 位通过该采样判据 \\
v3 reconstruct & status=0x00008026；固件输出 PASS & 同次上电注册后，重建状态未报告 mismatch/\allowbreak{}fallback \\
\bottomrule
\end{longtable}
}

固件 \texttt{keys\_\allowbreak{}rv32/\allowbreak{}main.\allowbreak{}c} 的 PASS 检查的是 ready、mismatch 与 fallback 状态，并未在这份板级日志中直接比较两次秘密 key，也没有以派生 key 完成 AEAD Tag 验证。因此``同次注册/重建自检通过''成立，``跨上电同密钥稳定''\,``认证启动通过''不成立。

\Needspace{85pt}

\paragraph{2.3.4.3 安全与实现边界}

当前 key\_mmio 可回读 REF、MASK、RAW 与 DNA，puf\_key 的派生上下文和算法也可见。因此仅有``key 寄存器不可直接读取''不足以证明软件无法重算派生值；本版应称为可观测的密钥生成实验原型，不能称为已经实现密钥保密隔离。无 helper 或匹配失败时的 DNA-only 回退也不提供秘密密钥保证。

已有扩展源集的时序审查记录 no\_clock=10944、未约束内部端点=27312，位于 RO 采样域。正 WNS 只覆盖当时已约束路径，不代表该扩展全域时序签核；这些数值不能套到无该子系统的 PYNQ 发布版。跨上电、温压、多器件唯一性及小时级稳定性未测，不是本次提交的阻塞条件，但报告不将其写成通过。

本原型作为已实施的探索展示；后续若开展可靠性或安全隔离改造，须独立立项、构建与验收，不替换已验证的三模式 PYNQ 演示版本。

\evid{证据 src/\allowbreak{}rtl/\allowbreak{}keys/\allowbreak{}puf\_\allowbreak{}key.\allowbreak{}v；src/\allowbreak{}rtl/\allowbreak{}keys/\allowbreak{}key\_\allowbreak{}mmio.\allowbreak{}v；src/\allowbreak{}firmware/\allowbreak{}keys\_\allowbreak{}rv32/\allowbreak{}main.\allowbreak{}c；data/\allowbreak{}ascon/\allowbreak{}results/\allowbreak{}board/\allowbreak{}uart\_\allowbreak{}20261027\_\allowbreak{}keys3.\allowbreak{}txt；docs/\allowbreak{}io\_\allowbreak{}timing\_\allowbreak{}signoff.\allowbreak{}md 的 RO 采样域记录}

\Needspace{85pt}

\section{软件系统介绍}

\Needspace{85pt}

\subsection{软件整体介绍}

软件分为 RV32 裸机固件、PS Python 调用层与开发机验证脚本。历史 PL-only 使用 ascon\_rv32、ascon\_hw\_rv32、ascon\_ise\_rv32；当前发布使用 pynq\_svc，在同一固件中按 mode 选择 SW/HW/ISE。固件由 COE 随 bitstream 初始化，不在运行时动态加载。

PS 侧由 AsconHost、MailboxBackend、PynqUartTransport 与 mailbox\_protocol.py 完成请求、身份、 deadline 与 epoch 管理。超时需显式 reset/recover，不隐式重发；单客户端串行，未实现跨进程锁。主机侧由 sim/ascon、sim/pynq、sim/board 的脚本分别执行标准验证、桥/服务仿真和回归。

\figpdf{2-14}{fig_2_11}{软件分层与调用关系}{PS Python、RV32 服务固件与验证脚本， 标明 AXI-Lite/UART 帧协议与 MMIO/定制指令边界。}

{\small
\renewcommand{\thetable}{2-19}
\begin{longtable}{@{}>{\raggedright\arraybackslash}p{\dimexpr 0.2926\linewidth-2\tabcolsep\relax}>{\raggedright\arraybackslash}p{\dimexpr 0.7074\linewidth-2\tabcolsep\relax}@{}}
\caption{设计平台与工具链}\label{tab:2-19}\\
\toprule
\textbf{项} & \textbf{当前发布版本与历史边界} \\
\midrule
\endfirsthead
\multicolumn{2}{@{}l}{\footnotesize 续表 2-19}\\
\toprule
\textbf{项} & \textbf{当前发布版本与历史边界} \\
\midrule
\endhead
器件与板卡 & 工程 part=\texttt{xc7z020clg400-\allowbreak{}1}；正点原子领航者 ZYNQ7020；完整芯片丝印/速度等级仍待现场取证 \\
当前时钟 & PS7 FCLK0 配置 100 MHz；独立频率测量 not\_\allowbreak{}measured \\
历史 PL-only 时钟 & 板载 50 MHz $\rightarrow$ clk\_\allowbreak{}wiz\_\allowbreak{}0 $\rightarrow$ 配置 100 MHz；不可替代当前 PS 时钟说明 \\
存储 & 128 KiB BRAM，COE 初始化；当前固件 pynq-svc-v1 \\
EDA & Vivado 2025.2 SW6299465 / IP6300035；top=pynq\_\allowbreak{}wrapper \\
当前 IP & processing\_\allowbreak{}system7 5.5、axi\_\allowbreak{}interconnect 2.1、axi\_\allowbreak{}bram\_\allowbreak{}ctrl 4.1、blk\_\allowbreak{}mem\_\allowbreak{}gen 8.4、proc\_\allowbreak{}sys\_\allowbreak{}reset 5.0；以 manifest 为准 \\
当前固件编译 & 归档构建使用 GCC 13.2.0，rv32im\_\allowbreak{}zicsr / ilp32、-O2；以构建日志和 ELF/BIN 哈希为准 \\
历史固件编译 & PL-only 基准使用 xpack GCC 15.2.0；不与当前发布版混用 \\
仿真 & 归档 PYNQ 回归为 Icarus/vvp 12.0 devel、Python 3.14.3；模块另有 XSim 2025.2 证据 \\
板端 & PYNQ 2.7.0、Python 3.8.2、Linux 5.\allowbreak{}4.\allowbreak{}0-\allowbreak{}xilinx-\allowbreak{}v2020.\allowbreak{}2 armv7l \\
\bottomrule
\end{longtable}
}

\evid{证据 board/\allowbreak{}vivado\_\allowbreak{}zynq7000\_\allowbreak{}xc7z020\_\allowbreak{}pynq/\allowbreak{}release/\allowbreak{}manifest.\allowbreak{}json；validation\_\allowbreak{}20261028.\allowbreak{}md；validation\_\allowbreak{}20261029.\allowbreak{}md；build/\allowbreak{}reports/\allowbreak{}pynq\_\allowbreak{}final\_\allowbreak{}20261028/\allowbreak{}}

\Needspace{85pt}

\subsection{软件各模块介绍}

\Needspace{85pt}

\paragraph{2.4.2.1 纯软件 AEAD 固件}\label{ux7eafux8f6fux4ef6-aead-ux56faux4ef6}

纯软件路径调用生产 ascon/ascon\_c/aead 的 crypto\_aead\_encrypt/decrypt，按 Final 参数完成初始化、AD 吸收、域分离、消息处理与终结。tag 采用完整 16 B。解密输入是密文与 tag 拼接，认证失败返回失败、输出长度归零并擦除候选明文。锁定上游参考不改；生产失败处理与原参考失败分别保留证据。

当前服务对 PT/AD 统一限制为 16 B；host 完整 KAT 与板级容量内 KAT 分列。固定 \#529 的当前 warm 加密中位为 23342 cycles，历史 PL-only 20806 cycles 仅作独立对照。

\Needspace{85pt}

\paragraph{2.4.2.2 MMIO 加速器驱动固件}\label{mmio-ux52a0ux901fux5668ux9a71ux52a8ux56faux4ef6}

MMIO 驱动使用 0x30000000 寄存器窗口：写 KEY(+0x10--1C)、NONCE(+0x20--2C)、 ADLEN/MLEN，推入 AD\_WDATA(+0x40)/MSG\_WDATA(+0x44)，写 CTRL 启动，轮询 STATUS，读取 OUT\_RDATA(+0x48) 与 TAGOUT(+0x4C--58)，结束后显式清零。消息与 AD 各 $\leq$16 B，BUSY 写忽略，认证成功且在有效长度内才允许输出。

当前 \#529 加密中位 1439 cycles，解密 1388 cycles；两者包含驱动调用内搬运与轮询，不含 PS/UART。 HW 轮询异常当前可能表现为 INTERNAL 或认证失败，不能声称存在独立 TIMEOUT 状态码；异常后整链恢复。历史 688 cycles 与核本体 103 cycles 不替代当前服务测量。

\Needspace{85pt}

\paragraph{2.4.2.3 定制指令固件}

ISE 用 ascon.w/ascon.rd 访问 10 个 32-bit 状态字，以 ascon.p12/ascon.p8 执行置换。主操作码 0x0B、funct7=0、funct3=0/1/2/3；inline assembly 使用 .insn，未修改 GCC backend。构建归档检出 24 条 custom-0 静态编码；真实服务运行 KAT 提供动态正确性证据。

当前 \#529 加密为 3995 cycles，相对 SW 为 5.843$\times$；历史 PL-only 2570 cycles/8.10$\times$ 单列。取消与 hold 的定向测试通过，不据此声称完整中断/抢占上下文保存恢复已验收。

\Needspace{85pt}

\paragraph{2.4.2.4 当前 benchmark 与采样工具}\label{ux5f53ux524d-benchmark-ux4e0eux91c7ux6837ux5de5ux5177}

board\_benchmark.py 对同一 release 完成身份检查、三模式 smoke 和固定输入采样。官方 KAT \#529、每模式预热 5 次，每操作 30 次，逐次核对 CT/tag/PT，并保存 raw\_cycles/raw\_wall\_ns。 JSON 同时记录服务标志、发布哈希、测量范围和统计；错误结果不能混入成功性能统计。

周期来自 RV32 mcycle 的高低高一致读取。PS 端墙钟独立计量，不以事件计数器相加代替总周期，也不将 CPU 周期加速比称为 PS 响应加速比。当前结果目录为 data/pynq/notebook\_20261029\_130920/。

\Needspace{85pt}

\paragraph{2.4.2.5 PS 端服务与帧协议}\label{ps-ux7aefux670dux52a1ux4e0eux5e27ux534fux8bae}

服务固件按 magic 识别 Y01E/Y01Q，校验 version、flags、epoch、request\_id、mode、operation 与长度，按 mode 调度三种实现，发送 Y01R；Y01S 发送固件身份、自检与 epoch 应答。当前帧没有 CRC 字段。请求/响应头各 48 B，Y01S 为 64 B，Y01E 为 20 B；失败响应无 payload/cycles。

超时不自动重传，须主机 reset/recover、更新 epoch 并重新验证 KAT。重复 ID 与错误 epoch 返回 INTERNAL；PT/AD \textgreater16 B 返回 LENGTH。半帧与 FIFO overflow 后恢复已板测；服务端无本地半帧接收超时，单客户端无跨进程锁。 9/9 错误门禁与每模式 50 次往返有 JSON/log/已执行 Notebook 证据，非小时级验证。

\Needspace{85pt}

\section{大模型协作与技能包}

\Needspace{85pt}

\subsection{协作分工与智能体工作流}

本作品按选题指南 3.3.3.1 使用大模型智能体辅助开发。分工原则是：\textbf{团队定目标、划边界、做裁决和实板验证，智能体在边界内执行编码、构建、仿真与资料整理}。凡是影响结论的判断，例如是否降频、报告声明什么范围、哪个版本可以上板，都由团队作出；智能体的输出只有经过脚本、仿真或实板复核后才进入成果。

{\small
\renewcommand{\thetable}{2-20}
\begin{longtable}{@{}>{\raggedright\arraybackslash}p{\dimexpr 0.1043\linewidth-2\tabcolsep\relax}>{\raggedright\arraybackslash}p{\dimexpr 0.3655\linewidth-2\tabcolsep\relax}>{\raggedright\arraybackslash}p{\dimexpr 0.2847\linewidth-2\tabcolsep\relax}>{\raggedright\arraybackslash}p{\dimexpr 0.2456\linewidth-2\tabcolsep\relax}@{}}
\caption{协作分工与智能体工作流}\label{tab:2-20}\\
\toprule
\textbf{环节} & \textbf{团队负责} & \textbf{智能体负责} & \textbf{交接物} \\
\midrule
\endfirsthead
\multicolumn{4}{@{}l}{\footnotesize 续表 2-20}\\
\toprule
\textbf{环节} & \textbf{团队负责} & \textbf{智能体负责} & \textbf{交接物} \\
\midrule
\endhead
目标与范围 & 确定 100 MHz 不降频、PT/AD $\leq$16 B 的交付范围，删去证据不足的功能声明 & 按范围列出缺口与工作量 & 任务看板、决策记录 \\
约束 & 维护仓库规则：不编造数据、未测写 not\_\allowbreak{}measured、不自动提交、不改生成物 & 每次会话先读规则再动手 & \texttt{AGENTS.\allowbreak{}md}、\texttt{docs/\allowbreak{}technical\_\allowbreak{}\allowbreak{}rules.\allowbreak{}md} \\
方案 & 审阅端点归因、选择改造顺序与候选方案 & 给出多套方案及代价，执行被选方案 & 方案文档、构建脚本 \\
执行 & 指定可改目录，安排多个智能体并行 & 父会话拆分子任务；子代理只改自己负责的目录 & 子代理会话、逐文件哈希 \\
复核 & 交叉审阅另一模型的结论，上板运行，判定是否通过 & 跑回归、汇总报告、记录失败 & 回归 JSON、板测日志 \\
沉淀 & 提出 Skill 要解决的问题与验证要求 & 整理成 Skill、测试与对照实验工具 & \texttt{skill/\allowbreak{}} 目录 \\
\bottomrule
\end{longtable}
}

智能体工作流采用''父会话规划、子代理分区执行、团队验收''。例如 10 月 29---30 日，团队要求''除上板外全部完成''后，父会话把任务拆为 PUF 数字部分、CPU 侧集成、独立审查等子代理，每个子代理只拥有写定的目录，并共享同一 Vivado 会话的使用约定；实板测试与是否采纳由团队决定。

\Needspace{85pt}

\subsection{代表性协作记录与纠错轨迹}

下列案例均摘自 \texttt{report/\allowbreak{}llm\_\allowbreak{}collab/\allowbreak{}sessions/\allowbreak{}} 中的原始会话，提示词为原文，模型回答为节选。每例给出团队提示、模型回答、纠错过程与最终裁决。

\textbf{案例一：拒绝降频，按端点归因完成 100 MHz 收敛（10-18）。}

{\small
\renewcommand{\thetable}{2-21}
\begin{longtable}{@{}>{\raggedright\arraybackslash}p{\dimexpr 0.1962\linewidth-2\tabcolsep\relax}>{\raggedright\arraybackslash}p{\dimexpr 0.8038\linewidth-2\tabcolsep\relax}@{}}
\caption{案例一：拒绝降频，按端点归因完成 100 MHz 收敛（10-18）}\label{tab:2-21}\\
\toprule
\textbf{步骤} & \textbf{记录} \\
\midrule
\endfirsthead
\multicolumn{2}{@{}l}{\footnotesize 续表 2-21}\\
\toprule
\textbf{步骤} & \textbf{记录} \\
\midrule
\endhead
团队提示 & "请你加载一下 AMD ROSS 的 skill$\ldots$$\ldots$先看看当前数据中的时序问题，并给我一个修改方案" \\
模型回答 & 策略扫掠后 WNS $-$0.177 ns；多次建议"降频到 98 MHz（强烈推荐）"，并给出"预期 WNS 0.523 ns"的样例输出 \\
团队裁决 & "不能降频，我们想完成时序修正" \\
纠错轨迹 & 两组实验 WNS 仍为 $-$0.177 ns，与基线完全相同。模型自查"是否真的应用了我们的修改"；团队判断"更改没有被识别到$\ldots$$\ldots$不可能进行了更改带来的 WNS 是一样的"，叫停模型，改为手动打开工程复核 \\
后续 & 团队在另一会话下达顺序：A（ISE 复位路径，预期清掉 93 个端点）$\rightarrow$ C（乘法器 48 个端点）$\rightarrow$ B（store/MMIO 绕行）$\rightarrow$ 双策略复测 \\
结果 & 两个策略在 100 MHz 下 WNS +0.001 / +0.007 ns，TNS 0；模型同时注明余量仅 1--7 ps，不能据此称上板通过（见 §3.3.1） \\
\bottomrule
\end{longtable}
}

\textbf{案例二：交叉审阅另一模型的结论，报告收缩到有证据的范围（10-29）。}

{\small
\renewcommand{\thetable}{2-22}
\begin{longtable}{@{}>{\raggedright\arraybackslash}p{\dimexpr 0.1962\linewidth-2\tabcolsep\relax}>{\raggedright\arraybackslash}p{\dimexpr 0.8038\linewidth-2\tabcolsep\relax}@{}}
\caption{案例二：交叉审阅另一模型的结论，报告收缩到有证据的范围（10-29）}\label{tab:2-22}\\
\toprule
\textbf{步骤} & \textbf{记录} \\
\midrule
\endfirsthead
\multicolumn{2}{@{}l}{\footnotesize 续表 2-22}\\
\toprule
\textbf{步骤} & \textbf{记录} \\
\midrule
\endhead
团队提示 & "这是另一个 ai 分析的，但是我感觉他的回复有可能是错的，请你给我分析一下现状距离我们想得到的还差什么" \\
团队裁决 & "先将 PUF$\ldots$$\ldots$、独立核 PPA 移除掉$\ldots$$\ldots$先保证我们的技术报告可以和当前的版本和证据对上" \\
纠错轨迹 & 按原始报告与 JSON 复算，纠正 23 类问题，如 MMIO/ISE 倍数写反、Vivado 版本写错、PUF key 位宽写成 320 bit、当前与历史周期混用；修订表见 \texttt{docs/\allowbreak{}tech\_\allowbreak{}\allowbreak{}report/\allowbreak{}review\_\allowbreak{}\allowbreak{}20261029.\allowbreak{}md} \\
结果 & 当前 PYNQ 数据与历史 PL-only 数据分列；证据脚本与一致性检查纳入一键构建 \\
\bottomrule
\end{longtable}
}

\textbf{案例三：团队发现板测漏接 PUF，重新划定测试组合（10-30）。}

{\small
\renewcommand{\thetable}{2-23}
\begin{longtable}{@{}>{\raggedright\arraybackslash}p{\dimexpr 0.1962\linewidth-2\tabcolsep\relax}>{\raggedright\arraybackslash}p{\dimexpr 0.8038\linewidth-2\tabcolsep\relax}@{}}
\caption{案例三：团队发现板测漏接 PUF，重新划定测试组合（10-30）}\label{tab:2-23}\\
\toprule
\textbf{步骤} & \textbf{记录} \\
\midrule
\endfirsthead
\multicolumn{2}{@{}l}{\footnotesize 续表 2-23}\\
\toprule
\textbf{步骤} & \textbf{记录} \\
\midrule
\endhead
团队提示 & "今天我们发现之前做的测试中 PUF 并没有接入$\ldots$$\ldots$你先给我说一下你的思路，我确认后，你进行操作" \\
模型回答 & 核对构建脚本后确认"PYNQ 上做过的所有测试里，PUF 确实没有接入"，并给出合并方案 \\
团队追问 & "这里的高速通道带来的优化有什么，如果$\ldots$$\ldots$需要分开测试，为什么我们要加入这个高速通道呢" \\
模型回答 & 分析指出高速通道绕开 RV32、与 PUF 无功能联系，合并只省一次下载却多一个未验证的 bit；建议不合并，给出 A/B/C 三案 \\
团队裁决 & 选定方案 A：PUF 第一版接入 PYNQ 新 overlay，高速通道用已有候选单独测，其他扩展后续单独修 \\
团队补充 & 指出另一版本的 setup.tcl 未把 PS 的 FIXED\_\allowbreak{}IO/DDR 引出、日志中有 Netlist 29-160 严重警告，判定该 bit 不得上板，并要求给出最小闭环 PUF 方案 \\
结果 & 当前报告只把 PUF 列为独立扩展原型（§2.3.4），不并入 PYNQ 主线成果 \\
\bottomrule
\end{longtable}
}

\textbf{案例四：版式迭代中的模型自我纠错（10-30）。}

团队逐页审阅 PDF，先后指出图被单独排成一页、公式与引导句分页、标题与图分离。模型第一次在每个标题前强制排出浮动图，虽消除了分离，却留下约 1/4--1/3 页空白；团队反馈''图离得太远''后，模型撤销该做法，改为按本页剩余版面决定原位、缩小或浮动，并用逐页填充率扫描确认除各部分末页外无大片空白。同一轮中，构建因 PDF 被阅读器占用而失败，模型依据非零退出码停止发布，待团队关闭文件后重跑全流程。

\Needspace{85pt}

\subsection{技能包的提炼与验证}

技能包位于 \texttt{skill/\allowbreak{}}，每项都从开发中的具体失败或观察提炼，并写明适用范围、使用方法、失效条件与已验证程度。

{\small
\renewcommand{\thetable}{2-24}
\begin{longtable}{@{}>{\raggedright\arraybackslash}p{\dimexpr 0.1378\linewidth-2\tabcolsep\relax}>{\raggedright\arraybackslash}p{\dimexpr 0.2874\linewidth-2\tabcolsep\relax}>{\raggedright\arraybackslash}p{\dimexpr 0.2874\linewidth-2\tabcolsep\relax}>{\raggedright\arraybackslash}p{\dimexpr 0.2874\linewidth-2\tabcolsep\relax}@{}}
\caption{技能包的提炼与验证}\label{tab:2-24}\\
\toprule
\textbf{Skill} & \textbf{来源（失败或观察）} & \textbf{内容} & \textbf{已验证程度} \\
\midrule
\endfirsthead
\multicolumn{4}{@{}l}{\footnotesize 续表 2-24}\\
\toprule
\textbf{Skill} & \textbf{来源（失败或观察）} & \textbf{内容} & \textbf{已验证程度} \\
\midrule
\endhead
\texttt{rtl-\allowbreak{}design-\allowbreak{}discipline} & 队友观察到 AI 写 Verilog 的通病：同一逻辑多处等价重写、打补丁式修改、不关心综合出的硬件；取指单元中同一选择表达式被嵌套重写三次 & 先列寄存器清单、每个信号只在一处赋值、互斥条件用 case、改需求时重写整段选择逻辑等可判定规则 & 文本已成形；用与不用的对照实验未做（not\_\allowbreak{}validated） \\
\texttt{pynq-\allowbreak{}overlay-\allowbreak{}workflow} & 本项目 PYNQ 实板中的真实问题：bit/hwh 不配套、ap\_\allowbreak{}done 读清、无界等待、参考时钟假设、缓存一致性 & Overlay 加载与哈希校验、寄存器驱动与测试骨架、黄金参考比对、性能与报告采集工具；示例不含 Ascon 地址或长度 & 工具层 47 项单元测试通过、13/13 变异检出；智能体对照与实板未运行 \\
\bottomrule
\end{longtable}
}

验证方法由团队提出要求后设计。团队指出：``加入 skill 和不加入 skill 在简单 RTL 生成时并无太大区别，并且在同一 agent 情况下去测试无法控制变量。''据此，\texttt{skill/\allowbreak{}ab\_\allowbreak{}eval\_\allowbreak{}pynq/\allowbreak{}} 采用三组受控对照：有 Skill、无 Skill、安慰剂 Skill（长度相同、同样会被加载，但只讲代码风格），每组使用新会话与独立工作目录，由判分脚本而不是模型自评打分，并先用''好答案/差答案''自检判分器能区分优劣。该实验尚未执行，报告不声称 Skill 已提升智能体效果。

边界的确定同样来自失败记录：Skill 不处理跨模块时序收敛（案例一说明这需要端点归因与 RTL 改造），也不替代实板验证；遇到这类问题时，Skill 的说明会引导使用者转向 Vivado 报告与板级测试。

\Needspace{85pt}

\subsection{原始记录的导出与范围}

原始记录由 \texttt{report/\allowbreak{}llm\_\allowbreak{}collab/\allowbreak{}export\_\allowbreak{}sessions.\allowbreak{}py} 从本机 OpenCode、Claude Code 与 Codex 的会话库只读导出，每个会话一份 Markdown，索引见 \texttt{report/\allowbreak{}llm\_\allowbreak{}collab/\allowbreak{}index.\allowbreak{}md}。导出时只做四项处理：工具输出折叠为一行摘要、超长回答保留首尾、疑似密钥打码，并省略 1 条与工程内容无关的用户提示原文（保留轮次与计数）；其余提示词与回答不改写。

截至 10 月 30 日的导出包含 21 个会话，其中 13 个为有工具调用的工程会话（76 条团队提示、约 2600 次工具调用），其余为无工具调用的问答和模型/工具自检，一并保留。所用模型包括 Claude Opus 系列与 GPT 系列，各会话的模型标识见索引。

记录存在三点范围限制：只含本机工具保存的会话；队友设备上的会话、网页对话及 10 月 9 日之前的工作未收录；同一结论若经多个会话迭代，以最终的工程证据（回归 JSON、实现报告、板测日志）为准，而不以模型回答为准。开发过程中的失败与返工另记录在 \texttt{project\_\allowbreak{}log/\allowbreak{}lessons\_\allowbreak{}learned.\allowbreak{}md}。
